\chapter{Transformasi Laplace} \label{LT:chapter}

%%%%%%%%%%%%%%%%%%%%%%%%%%%%%%%%%%%%%%%%%%%%%%%%%%%%%%%%%%%%%%%%%%%%%%%%%%%%%%

\section{Transformasi Laplace}
\label{laplace:section}

%mbxINTROSUBSECTION

\sectionnotes{1.5--2 kuliah\EPref{, \S7.1 dalam \cite{EP}}\BDref{,
\S6.1 dan bagian-bagian dari \S6.2 dalam \cite{BD}}}

\subsection{Transformasi}

Dalam bab ini, kita akan membahas transformasi Laplace%
\footnote{Sama seperti persamaan Laplace dan operator Laplace, transformasi
Laplace juga dinamai menurut 
\href{https://en.wikipedia.org/wiki/Laplace}{Pierre-Simon, marquis de Laplace}
(1749--1827).}.
Transformasi Laplace
adalah metode yang sangat efisien untuk menyelesaikan masalah ODE atau PDE tertentu.
Transformasi tersebut mengubah persamaan diferensial menjadi
persamaan aljabar.  Jika persamaan aljabar itu dapat diselesaikan, penerapan
transformasi invers memberi kita solusi yang diinginkan.
Transformasi Laplace juga memiliki penerapan dalam
analisis 
rangkaian listrik, spektroskopi NMR, pemrosesan sinyal, dan bidang lainnya.
Terakhir,
memahami transformasi
Laplace juga akan membantu memahami transformasi Fourier yang berkaitan,
yang, bagaimanapun, memerlukan pemahaman lebih mendalam
tentang bilangan kompleks.  Kita tidak akan membahas transformasi Fourier.

Transformasi Laplace juga memberikan banyak wawasan tentang sifat
persamaan yang kita hadapi.  Transformasi ini dapat dipandang sebagai pengubahan antara ranah
waktu dan ranah frekuensi.  Sebagai contoh, ambil persamaan standar
\begin{equation*}
m x''(t) + c x'(t) + k x(t) = f(t) .
\end{equation*}
Kita dapat memandang $t$ sebagai waktu dan $f(t)$ sebagai sinyal masuk.  Transformasi Laplace
akan mengubah persamaan tersebut dari persamaan diferensial dalam waktu menjadi
persamaan aljabar (tanpa turunan), dengan peubah bebas baru
$s$ yang dipandang sebagai frekuensi.\footnote{Sesungguhnya, ini adalah \myquote{frekuensi} dalam konteks
bilangan kompleks, tetapi kita menyimpang dari pokok bahasan.}

Kita dapat memandang \emph{\myindex{transformasi Laplace}} sebagai kotak hitam.  Kotak ini
menerima fungsi dan menghasilkan fungsi dalam peubah baru.  Kita menulis
$\mathcal{L} \bigl\{ f(t) \bigr\} = F(s)$ untuk transformasi Laplace dari $f(t)$.
Biasanya huruf kecil digunakan untuk
fungsi dalam ranah waktu dan huruf besar untuk fungsi dalam
ranah frekuensi.  Kita menggunakan huruf yang sama untuk menunjukkan bahwa satu fungsi
merupakan transformasi Laplace dari fungsi lainnya.  Sebagai contoh, $F(s)$ adalah transformasi Laplace
dari $f(t)$.  Mari kita definisikan transformasi tersebut.
\begin{equation*}
\mybxbg{~~
\mathcal{L} \bigl\{ f(t) \bigr\} =
F(s) \overset{\text{def}}{=} \int_0^\infty e^{-st} f(t) \, dt .
~~}
\end{equation*}
Perhatikan bahwa kita hanya mempertimbangkan $t \geq 0$ dalam transformasi tersebut.  Tentu saja,
jika kita memandang $t$ sebagai waktu, tidak ada masalah; pada umumnya kita tertarik untuk
mengetahui apa yang akan terjadi pada masa depan
(transformasi Laplace adalah salah satu tempat
yang aman untuk mengabaikan masa lalu).
Mari kita hitung beberapa
transformasi sederhana.

\begin{example}
Andaikan $f(t) = 1$, maka
\begin{equation*}
\mathcal{L} \{1\} = \int_0^\infty e^{-st} \, dt
=
\left[ \frac{e^{-st}}{-s} \right]_{t=0}^\infty
=
\lim_{h\to\infty}
\left[ \frac{e^{-st}}{-s} \right]_{t=0}^h
=
\lim_{h\to\infty}
\left( \frac{e^{-sh}}{-s} - \frac{1}{-s} \right)
= \frac{1}{s} .
\end{equation*}
Limit tersebut (integral takwajar) hanya ada jika $s > 0$.  Jadi, 
$\mathcal{L} \{1\}$ hanya terdefinisi untuk $s > 0$.
\end{example}

\begin{example}
Andaikan $f(t) = e^{-at}$, maka
\begin{equation*}
\mathcal{L} \bigl\{e^{-at}\bigr\}
= \int_0^\infty e^{-st} e^{-at} \, dt
= \int_0^\infty e^{-(s+a)t} \, dt
=
\left[ \frac{e^{-(s+a)t}}{-(s+a)} \right]_{t=0}^\infty
= \frac{1}{s+a} .
\end{equation*}
Limit tersebut hanya ada jika $s+a > 0$.  Jadi, 
$\mathcal{L} \bigl\{e^{-at}\bigr\}$ hanya terdefinisi untuk $s+a > 0$.
\end{example}

\begin{example}
Andaikan $f(t) = t$, maka dengan menggunakan integrasi parsial
\begin{equation*}
\begin{split}
\mathcal{L} \{t\}
& = \int_0^\infty e^{-st} t \, dt \\
& =
\left[ \frac{-te^{-st}}{s} \right]_{t=0}^\infty
+
\frac{1}{s}
\int_0^\infty e^{-st} \,dt \\
& =
0
+
\frac{1}{s}
\left[ \frac{e^{-st}}{-s} \right]_{t=0}^\infty \\
& =
\frac{1}{s^2} .
\end{split}
\end{equation*}
Sekali lagi, limit tersebut hanya ada jika $s > 0$.
\end{example}

\begin{example}
Salah satu fungsi yang umum adalah \emph{\myindex{fungsi tangga satuan}}, yang
kadang-kadang disebut \emph{\myindex{fungsi Heaviside}}%
\footnote{Fungsi tersebut dinamai menurut matematikawan, insinyur, dan
fisikawan Inggris
\href{https://en.wikipedia.org/wiki/Heaviside}{Oliver Heaviside}
(1850--1925).  Hanya
secara kebetulan fungsi ini \myquote{berat} di \myquote{satu sisi}.}.
Fungsi ini umumnya diberikan sebagai
\begin{equation*}
u(t) = 
\begin{cases}
0 & \text{jika } \; t < 0 , \\
1 & \text{jika } \; t \geq 0 .
\end{cases}
\end{equation*}
%Note that some authors prefer to define $u(0) = \frac{1}{2}$.
Mari kita cari transformasi Laplace dari $u(t-a)$, dengan $a \geq 0$
suatu konstanta.
Artinya, ini adalah fungsi yang bernilai $0$ untuk $t < a$ dan $1$ untuk $t \geq a$.
\begin{equation*}
\mathcal{L} \bigl\{ u(t-a) \bigr\}
=
\int_0^{\infty} e^{-st} u(t-a) \, dt
=
\int_a^{\infty} e^{-st} \, dt
=
\left[ \frac{e^{-st}}{-s} \right]_{t=a}^\infty
=
\frac{e^{-as}}{s} ,
\end{equation*}
dengan, tentu saja, $s > 0$ (dan $a \geq 0$ seperti yang telah kita nyatakan).
\end{example}

Dengan menerapkan prosedur serupa, kita dapat menghitung transformasi dari banyak
fungsi elementer.  Banyak transformasi dasar dicantumkan dalam
\tablevref{lt:table} (lihat juga \appendixref{laplacelist:appendix}).

\begin{table}[h!t]
\mybeginframe
\capstart
\begin{center}
\begin{tabular}{@{}lllll@{}}
\toprule
$f(t)$ & $\mathcal{L} \bigl\{ f(t) \bigr\}$ & $\qquad\quad$ &
$f(t)$ & $\mathcal{L} \bigl\{ f(t) \bigr\}$ \\
\midrule
$C$ & $\frac{C}{s}$
& &
$\sin (\omega t)$ & $\frac{\omega}{s^2+\omega^2}$
\\[4pt]
$t$ & $\frac{1}{s^2}$
& &
$\cos (\omega t)$ & $\frac{s}{s^2+\omega^2}$
\\[4pt]
$t^2$ & $\frac{2}{s^3}$
& &
$\sinh (\omega t)$ & $\frac{\omega}{s^2-\omega^2}$
\\[4pt]
$t^3$ & $\frac{6}{s^4}$
& &
$\cosh (\omega t)$ & $\frac{s}{s^2-\omega^2}$
\\[4pt]
$t^n$ & $\frac{n!}{s^{n+1}}$
& &
$u(t-a) \quad (a \geq 0)$ & $\frac{e^{-as}}{s}$
\\[4pt]
$e^{-at}$ & $\frac{1}{s+a}$
& & &
\\[4pt]
\bottomrule
\end{tabular}
\end{center}
\caption{Beberapa transformasi Laplace ($C$, $\omega$, dan $a$ adalah
konstanta).\label{lt:table}}
\myendframe
\end{table}

\begin{exercise}
Verifikasikan \tablevref{lt:table}.
\end{exercise}

Karena transformasi didefinisikan oleh integral, kita dapat menggunakan sifat
linearitas integral.  Sebagai contoh, jika $C$ adalah konstanta, maka
\begin{equation*}
\mathcal{L} \bigl\{ C f(t) \bigr\} =
\int_0^\infty e^{-st} C f(t) \,dt =
C \int_0^\infty e^{-st} f(t) \,dt =
C \mathcal{L} \bigl\{ f(t) \bigr\} .
\end{equation*}
Jadi, kita dapat \myquote{mengeluarkan} konstanta dari transformasi.
Demikian pula untuk penjumlahan.
Karena linearitas penting, kita menyatakannya sebagai teorema.

\begin{theorem}[Linearitas transformasi Laplace]
\index{linearitas transformasi Laplace}
Jika $A$, $B$, dan $C$ adalah konstanta, maka
\begin{equation*}
\mybxbg{~~
\mathcal{L} \bigl\{ A f(t) + B g(t) \bigr\} =
A \mathcal{L} \bigl\{ f(t) \bigr\} +
B \mathcal{L} \bigl\{ g(t) \bigr\} ,
~~}
\end{equation*}
dan khususnya,
\begin{equation*}
\mathcal{L} \bigl\{ C f(t) \bigr\} =
C \mathcal{L} \bigl\{ f(t) \bigr\} .
\end{equation*}
\end{theorem}

\begin{exercise}
Verifikasikan teorema tersebut.  Artinya, tunjukkan bahwa
$\mathcal{L} \bigl\{ A f(t) + B g(t) \bigr\} =
A \mathcal{L} \bigl\{ f(t) \bigr\} +
B \mathcal{L} \bigl\{ g(t) \bigr\}$.
\end{exercise}

Aturan-aturan ini bersama dengan \tablevref{lt:table} sudah memudahkan kita mencari
transformasi Laplace dari sangat banyak fungsi.
Sebagai contoh:
\begin{equation*}
\mathcal{L}\bigl\{
2 + 5t + 9e^{-2t}
\bigr\}
=
\mathcal{L}\bigl\{
2
\bigr\}
+5
\mathcal{L}\bigl\{
t
\bigr\}
+
9
\mathcal{L}\bigl\{
e^{-2t}
\bigr\}
=
\frac{2}{s} + \frac{5}{s^2} + \frac{9}{s+2} .
\end{equation*}
Berhati-hatilah!
Transformasi Laplace dari suatu hasil kali
\emph{bukan} hasil kali transformasi-transformasinya.  Secara umum, 
\begin{equation*}
\mathcal{L} \bigl\{ f(t) g(t) \bigr\} \neq
\mathcal{L} \bigl\{ f(t) \bigr\}
\mathcal{L} \bigl\{ g(t) \bigr\} .
\end{equation*}
Selain itu, tidak semua fungsi memiliki transformasi Laplace.  Sebagai
contoh, fungsi $\frac{1}{t}$ tidak memiliki transformasi Laplace karena
integralnya divergen untuk semua $s$.  Demikian pula,
$\tan t$ atau $e^{t^2}$ tidak memiliki transformasi Laplace.

\subsection{Eksistensi dan ketunggalan}

Kapan transformasi Laplace ada?  Suatu fungsi $f(t)$ dikatakan ber-
\emph{\myindex{orde eksponensial}} ketika $t$ menuju takhingga jika
\begin{equation*}
\lvert f(t) \rvert \leq M e^{ct} ,
\end{equation*}
untuk suatu konstanta $M$ dan $c$, untuk
$t$ yang cukup besar (katakanlah untuk semua $t > t_0$ bagi suatu $t_0$).  Cara termudah
untuk memeriksa syarat ini adalah mencoba menghitung
\begin{equation*}
\lim_{t\to \infty} \frac{f(t)}{e^{ct}} .
\end{equation*}
Jika limit tersebut ada dan berhingga (biasanya nol), maka $f(t)$ ber-
orde eksponensial.

\begin{exercise}
Gunakan aturan L'Hopital dari kalkulus untuk menunjukkan bahwa suatu polinom ber-
orde eksponensial.  Petunjuk: Perhatikan bahwa jumlah dua fungsi berorde eksponensial
juga berorde eksponensial.  Kemudian tunjukkan bahwa $t^n$ berorde eksponensial
untuk setiap $n$.
\end{exercise}

Untuk fungsi berorde eksponensial, kita memiliki eksistensi dan ketunggalan
transformasi Laplace.

\begin{theorem}[Eksistensi]
Misalkan $f(t)$ kontinu pada interval $[0,\infty)$ dan berorde eksponensial untuk suatu
konstanta $c$.  Maka $F(s) = \mathcal{L} \bigl\{ f(t) \bigr\}$ terdefinisi untuk
semua $s > c$.
\end{theorem}

Eksistensinya tidak sulit dilihat.  Misalkan $f(t)$ berorde eksponensial,
yakni, $\lvert f(t) \rvert \leq M e^{ct}$ untuk semua $t > 0$ (demi kesederhanaan $t_0 = 0$).
Misalkan $s > c$, atau dengan kata lain $(s-c) > 0$.
Berdasarkan teorema perbandingan dari kalkulus, integral takwajar yang mendefinisikan
$\mathcal{L} \bigl\{ f(t) \bigr\}$ ada karena integral berikut ada:
\begin{equation*}
\int_0^\infty e^{-st} ( M e^{ct} ) \,dt
=
M \int_0^\infty e^{-(s-c)t} \,dt = M \left[ \frac{e^{-(s-c)t}}{-(s-c)}
\right]_{t=0}^\infty = \frac{M}{s-c} .
\end{equation*}

Transformasi tersebut juga ada untuk beberapa fungsi lain
yang tidak berorde
eksponensial, tetapi hal itu tidak relevan bagi kita.
Sebelum membahas ketunggalan,
perhatikan bahwa untuk fungsi berorde eksponensial,
transformasi Laplace-nya meluruh di takhingga:
\begin{equation*}
\lim_{s\to\infty} F(s) = 0 .
\end{equation*}

\begin{theorem}[Ketunggalan] \label{lt:uniqthm}
Misalkan $f(t)$ dan $g(t)$ kontinu dan berorde eksponensial.
Andaikan terdapat suatu konstanta $C$
sedemikian sehingga $F(s) = G(s)$ untuk semua $s > C$.
Maka $f(t) = g(t)$ untuk semua $t \geq 0$.
\end{theorem}

Kedua teorema tersebut juga berlaku untuk fungsi kontinu sepotong-sepotong.
Ingat bahwa kontinu sepotong-sepotong berarti fungsi tersebut
kontinu kecuali mungkin pada suatu himpunan titik diskret, tempat ia memiliki
diskontinuitas loncatan seperti fungsi Heaviside.  Namun, transformasi Laplace tidak
\myquote{melihat} nilai-nilai di titik diskontinuitas.  Jadi, kita hanya dapat menyimpulkan bahwa
$f(t) = g(t)$ di luar titik-titik diskontinuitas.  Sebagai contoh, fungsi tangga
satuan kadang-kadang didefinisikan menggunakan $u(0) = \nicefrac{1}{2}$.  Namun, fungsi
tangga yang baru ini memiliki transformasi Laplace yang persis sama
dengan fungsi yang kita definisikan sebelumnya, dengan $u(0) = 1$.

\subsection{Transformasi invers}

Seperti yang telah kita katakan, transformasi Laplace memungkinkan kita mengubah persamaan
diferensial menjadi persamaan aljabar.  Setelah kita menyelesaikan
persamaan aljabar dalam ranah frekuensi, kita ingin kembali ke
ranah waktu karena itulah yang kita minati.
Diberikan suatu fungsi
$F(s)$, kita ingin mencari fungsi
$f(t)$ sedemikian sehingga $\mathcal{L} \bigl\{ f(t) \bigr\} = F(s)$.
\thmref{lt:uniqthm} menyatakan bahwa solusi $f(t)$ tunggal.
Jadi, tanpa kekhawatiran kita dapat membuat definisi berikut.

Andaikan $F(s) = \mathcal{L} \bigl\{ f(t) \bigr\}$ untuk suatu fungsi $f(t)$.
Definisikan
\emph{\myindex{transformasi Laplace invers}} sebagai
\begin{equation*}
{\mathcal{L}}^{-1} \bigl\{ F(s) \bigr\} \overset{\text{def}}{=} f(t) .
\end{equation*}
Terdapat rumus integral untuk invers tersebut, tetapi rumusnya tidak sesederhana
transformasi itu sendiri---rumus tersebut memerlukan bilangan kompleks dan integral lintasan.
Bagi kita, cukup
menghitung invers
menggunakan \tablevref{lt:table}.

\begin{example}
Ambil
$F(s) = \frac{1}{s+1}$.  Carilah transformasi Laplace inversnya.

Kita melihat tabel untuk memperoleh
\begin{equation*}
{\mathcal{L}}^{-1} \left\{ \frac{1}{s+1} \right\} = 
e^{-t} .
\end{equation*}
\end{example}

Karena transformasi Laplace linear, transformasi Laplace
invers juga linear.  Artinya,
\begin{equation*}
{\mathcal{L}}^{-1} \bigl\{ A F(s) + B G(s) \bigr\} =
A {\mathcal{L}}^{-1} \bigl\{ F(s) \bigr\} +
B {\mathcal{L}}^{-1} \bigl\{ G(s) \bigr\} .
\end{equation*}
Tentu saja, kita juga memiliki
${\mathcal{L}}^{-1} \bigl\{ A F(s) \bigr\} =
A {\mathcal{L}}^{-1} \bigl\{ F(s) \bigr\}$.
%Let us demonstrate how linearity can be used.

\begin{example}
Ambil
$F(s) = \frac{s^2+s+1}{s^3+s}$.  Carilah transformasi Laplace inversnya.

Pertama, kita menggunakan \emph{\myindex{metode pecahan parsial}} untuk menulis $F$ dalam bentuk yang
memungkinkan kita menggunakan \tablevref{lt:table}.  Kita faktorkan penyebut sebagai
$s(s^2+1)$ dan tulis
\begin{equation*}
\frac{s^2+s+1}{s^3+s}
=
\frac{A}{s} + 
\frac{Bs+C}{s^2+1} .
\end{equation*}
Dengan menyamakan penyebut pada ruas kanan
dan menyamakan pembilang, kita memperoleh
$A(s^2+1) + s(Bs+C) = s^2+s+1$.  Dengan menguraikan dan menyamakan koefisien,
kita memperoleh $A+B = 1$, $C=1$, $A=1$,
dan dengan demikian $B=0$.  Dengan
kata lain,
\begin{equation*}
F(s) =
\frac{s^2+s+1}{s^3+s}
=
\frac{1}{s} +
\frac{1}{s^2+1} .
\end{equation*}
Berdasarkan linearitas transformasi Laplace invers, kita memperoleh 
\begin{equation*}
{\mathcal{L}}^{-1} \left\{ 
\frac{s^2+s+1}{s^3+s} \right\}
=
{\mathcal{L}}^{-1} \left\{ 
\frac{1}{s} \right\} 
+
{\mathcal{L}}^{-1} \left\{ 
\frac{1}{s^2+1} \right\}
=
1 + 
\sin t .
\end{equation*}
\end{example}

Sifat berguna lainnya adalah apa yang disebut 
\emph{\myindex{sifat pergeseran}} atau
\emph{\myindex{sifat pergeseran pertama}}
\begin{equation*}
\mybxbg{~~
\mathcal{L} \bigl\{ e^{-at} f(t) \bigr\} = F(s+a) ,
~~}
\end{equation*}
dengan $F(s)$ transformasi Laplace dari $f(t)$.

\begin{exercise}
Turunkan sifat pergeseran pertama
dari definisi transformasi Laplace.
\end{exercise}

Sifat pergeseran dapat digunakan, misalnya, ketika penyebutnya berupa
bentuk kuadratik yang lebih rumit yang mungkin muncul dalam metode pecahan
parsial.  Kita melengkapkan kuadrat dan menulis bentuk kuadratik tersebut sebagai ${(s+a)}^2+b$,
lalu menggunakan sifat pergeseran.

\begin{example}
Carilah
${\mathcal{L}}^{-1} \left\{ \frac{1}{s^2+4s+8} \right\}$.

Pertama, kita melengkapkan kuadrat sehingga penyebut menjadi ${(s+2)}^2+4$.  
Kita memperoleh
\begin{equation*}
{\mathcal{L}}^{-1} \left\{ \frac{1}{s^2+4} \right\}
=
\frac{1}{2}
{\mathcal{L}}^{-1} \left\{ \frac{2}{s^2+2^2} \right\}
=
\frac{1}{2} \sin (2t) .
\end{equation*}
Terakhir, kita gabungkan semuanya dengan sifat pergeseran untuk memperoleh
\begin{equation*}
{\mathcal{L}}^{-1} \left\{ \frac{1}{s^2+4s+8} \right\} = 
{\mathcal{L}}^{-1} \left\{ \frac{1}{{(s+2)}^2+4} \right\}
=
\frac{1}{2}\,e^{-2t} \sin (2t) .
\end{equation*}
\end{example}

Kita sering ingin dapat menerapkan transformasi Laplace invers pada
fungsi rasional, yakni, fungsi berbentuk
\begin{equation*}
\frac{F(s)}{G(s)}
\end{equation*}
dengan $F(s)$ dan $G(s)$ polinom.  Jika $\frac{F(s)}{G(s)}$ adalah
transformasi Laplace dari fungsi berorde eksponensial,
nilainya menuju nol ketika $s \to \infty$, sehingga derajat $F(s)$
harus lebih kecil daripada derajat $G(s)$.  Fungsi rasional semacam itu
disebut \emph{fungsi rasional sejati\index{fungsi rasional sejati}},
dan kita selalu dapat menerapkan metode pecahan parsial tanpa pembagian
polinom.
Tentu saja, kita tetap perlu dapat memfaktorkan penyebut menjadi
suku-suku linear dan kuadratik, yang melibatkan pencarian akar
penyebut tersebut.

\subsection{Latihan}

\begin{exercise}
Carilah transformasi Laplace dari $3+t^5+\sin (\pi t)$.
\end{exercise}

\begin{exercise}
Carilah transformasi Laplace dari $a+bt+ct^2$ untuk suatu konstanta $a$, $b$, dan
$c$.
\end{exercise}

\begin{exercise}
Carilah transformasi Laplace dari $A \cos (\omega t) + B \sin (\omega t)$.
\end{exercise}

\begin{exercise}
Carilah transformasi Laplace dari $\cos^2 (\omega t)$.
\end{exercise}

\begin{exercise}
Carilah transformasi Laplace invers dari $\frac{4}{s^2-9}$.
\end{exercise}

\begin{exercise}
Carilah transformasi Laplace invers dari $\frac{2s}{s^2-1}$.
\end{exercise}

\begin{exercise}
Carilah transformasi Laplace invers dari $\frac{1}{{(s-1)}^2(s+1)}$.
\end{exercise}

\begin{exercise}
Carilah transformasi Laplace dari $f(t) =
\begin{cases}
t & \text{jika } \; t \geq 1, \\
0 & \text{jika } \; t < 1.
\end{cases}$
\end{exercise}

\begin{exercise}
Carilah transformasi Laplace invers dari $\frac{s}{(s^2+s+2)(s+4)}$.
\end{exercise}

\begin{exercise}
Carilah transformasi Laplace dari $\sin\bigl(\omega (t-a)\bigr)$.
\end{exercise}

\begin{exercise}
Carilah transformasi Laplace dari $t\sin(\omega t)$.  Petunjuk: Lakukan integrasi
parsial beberapa kali.
\end{exercise}

\setcounter{exercise}{100}

\begin{exercise}
Carilah transformasi Laplace dari $4{(t+1)}^2$.
\end{exercise}
\exsol{%
$\frac{8}{s^3} + \frac{8}{s^2} + \frac{4}{s}$
}

\begin{exercise}
Carilah transformasi Laplace invers dari $\frac{8}{s^3(s+2)}$.
\end{exercise}
\exsol{%
$2t^2-2t+1-e^{-2t}$
}

\begin{exercise}
Carilah transformasi Laplace dari $te^{-t}$.  Petunjuk: Sifat pergeseran.
\end{exercise}
\exsol{%
$\frac{1}{{(s+1)}^2}$
}

\begin{exercise}
Carilah transformasi Laplace dari $\sin(t)e^{-t}$.  Petunjuk: Integrasikan secara parsial.
\end{exercise}
\exsol{%
$\frac{1}{s^2+2s+2}$
}


%%%%%%%%%%%%%%%%%%%%%%%%%%%%%%%%%%%%%%%%%%%%%%%%%%%%%%%%%%%%%%%%%%%%%%%%%%%%%%

\sectionnewpage
\section{Transformasi turunan dan ODE}
\label{transformsofders:section}

%mbxINTROSUBSECTION

\sectionnotes{2 kuliah\EPref{, \S7.2--7.3 dalam \cite{EP}}\BDref{,
\S6.2 dan \S6.3 dalam \cite{BD}}}

\subsection{Transformasi turunan}

Mari kita lihat bagaimana transformasi Laplace digunakan untuk persamaan diferensial.
Pertama, kita cari transformasi Laplace dari turunan suatu fungsi.
Andaikan $g(t)$ adalah fungsi terdiferensialkan
berorde eksponensial, yakni, $\lvert g(t) \rvert \leq M e^{ct}$ untuk suatu
$M$ dan $c$.  Jadi, $\mathcal{L} \bigl\{ g(t) \bigr\}$ ada, dan terlebih lagi,
$\lim_{t\to\infty} e^{-st}g(t) = 0$ ketika $s > c$.  Maka
\begin{equation*}
\mathcal{L} \bigl\{ g'(t) \bigr\}
=
\int_0^\infty
e^{-st}
g'(t) \,dt
=
\Bigl[e^{-st} g(t) \Bigr]_{t=0}^\infty
-
\int_0^\infty
(-s)\,
e^{-st}
g(t) \,dt
=
-g(0) + s \mathcal{L} \bigl\{ g(t) \bigr\} .
\end{equation*}
Kita mengulangi prosedur ini untuk turunan tingkat lebih tinggi.
Hasilnya
dicantumkan dalam \tablevref{ltd:table}.  Prosedur tersebut juga berlaku untuk fungsi kontinu mulus
sepotong-sepotong, yakni, fungsi kontinu dengan turunan
kontinu sepotong-sepotong.
%The fact that the function is of
%exponential order is used to show that the limits appearing above 
%exist.  We will not worry much about this fact.

\begin{table}[h!t]
\mybeginframe
\capstart
\begin{center}
\begin{tabular}{@{}ll@{}}
\toprule
$f(t)$ & $\mathcal{L} \bigl\{ f(t) \bigr\} = F(s)$ \\
\midrule
$g'(t)$ & $sG(s)-g(0)$ \\[4pt]
$g''(t)$ & $s^2G(s)-sg(0)-g'(0)$ \\[4pt]
$g'''(t)$ & $s^3G(s)-s^2g(0)-sg'(0)-g''(0)$ \\[4pt]
\bottomrule
\end{tabular}
\end{center}
\caption{Transformasi Laplace dari turunan
($G(s) = \mathcal{L} \bigl\{ g(t) \bigr\}$
seperti biasa).\label{ltd:table}}
\myendframe
\end{table}

\begin{exercise}
Verifikasikan \tablevref{ltd:table}.
\end{exercise}

\subsection{Menyelesaikan ODE dengan transformasi Laplace}

Perhatikan bahwa transformasi Laplace mengubah pendiferensialan menjadi
perkalian dengan $s$.
Sifat inilah yang menjadikan transformasi tersebut berguna untuk diterapkan 
pada persamaan diferensial.

\begin{example}
Perhatikan masalah
\begin{equation*}
x''(t) + x(t) = \cos (2t), \quad x(0) = 0, \quad x'(0) = 1 .
\end{equation*}
Kita akan mengambil transformasi Laplace dari kedua ruas persamaan.
Seperti biasa, dengan $X(s)$ kita menyatakan transformasi Laplace dari
$x(t)$.
\begin{align*}
\mathcal{L} \bigl\{ x''(t) + x(t) \bigr\} & = \mathcal{L} \bigl\{ \cos (2t) \bigr\} , \\
s^2 X(s) -sx(0)-x'(0) + X(s) & = \frac{s}{s^2 + 4} .
\end{align*}
Sekarang kita substitusikan syarat awal---sehingga perhitungan menjadi lebih
ringkas---untuk memperoleh
\begin{equation*}
s^2 X(s) -1 + X(s) = \frac{s}{s^2 + 4} .
\end{equation*}
Kita selesaikan untuk $X(s)$,
\begin{equation*}
X(s) = \frac{s}{(s^2+1)(s^2 + 4)} + \frac{1}{s^2+1} .
\end{equation*}
Kita menggunakan pecahan parsial (latihan) untuk menulis
\begin{equation*}
X(s) =\frac{1}{3} \, \frac{s}{s^2+1} - 
\frac{1}{3}\, \frac{s}{s^2+4} + \frac{1}{s^2+1} .
\end{equation*}
Sekarang ambil transformasi Laplace invers untuk memperoleh
\begin{equation*}
x(t) =\frac{1}{3}  \cos (t) -
\frac{1}{3} \cos (2t) + \sin (t) .
\end{equation*}
\end{example}

Prosedur untuk persamaan linear berkoefisien konstan adalah sebagai berikut:
Ambil suatu persamaan diferensial biasa
dalam peubah waktu $t$.  Terapkan transformasi Laplace
untuk mengubah persamaan tersebut menjadi persamaan aljabar (bukan diferensial) dalam
ranah frekuensi.  Semua $x(t)$, $x'(t)$, $x''(t)$, dan seterusnya akan
diubah menjadi $X(s)$, $sX(s) - x(0)$, $s^2X(s) - sx(0) - x'(0)$,
dan seterusnya.
Selesaikan persamaan tersebut untuk $X(s)$.
Kemudian, dengan mengambil transformasi invers, jika memungkinkan, carilah $x(t)$.

Perlu diperhatikan bahwa karena tidak setiap fungsi memiliki transformasi Laplace,
tidak setiap persamaan dapat diselesaikan dengan cara ini.  Selain itu, jika persamaan
tersebut bukan ODE linear berkoefisien konstan\@,
maka penerapan transformasi Laplace mungkin tidak
menghasilkan persamaan aljabar.

\subsection{Menggunakan fungsi Heaviside}

Sebelum beralih ke persamaan yang lebih umum
daripada yang dapat kita selesaikan sebelumnya,
kita ingin membahas fungsi Heaviside.  Lihat \figurevref{lt:heavisidefig}
untuk grafiknya.
\begin{equation*}
u(t) =
\begin{cases}
0 & \text{jika } \; t < 0 , \\ 
1 & \text{jika } \; t \geq 0 .
\end{cases}
\end{equation*}

\begin{myfig}
\capstart
\diffyincludegraphics{width=3in}{width=4.5in}{lt-heaviside}{%
Grafik fungsi tangga satuan pada interval dari minus 1 hingga 1.
Fungsi bernilai nol untuk t negatif dan 1 untuk t positif.}
\caption{Grafik fungsi Heaviside (tangga satuan)
$u(t)$.\label{lt:heavisidefig}}
\end{myfig}

Fungsi Heaviside berguna untuk
menggabungkan fungsi lain atau memotong fungsi.  Biasanya kita menggunakan
$u(t-a)$ untuk suatu konstanta $a$;  kita hanya menggeser grafik ke
kanan sebesar $a$.  Artinya, $u(t-a) = 0$ ketika $t < a$ dan $u(t-a) = 1$
ketika $t \geq a$.  Sebagai contoh, andaikan $f(t)$ adalah suatu \myquote{sinyal},
dan kita mulai menerima $\sin t$ pada waktu $t=\pi$.
Fungsi $f(t)$ kemudian didefinisikan sebagai
\begin{equation*}
f(t) =
\begin{cases}
0 & \text{jika } \; t < \pi , \\ 
\sin t & \text{jika } \; t \geq \pi .
\end{cases}
\end{equation*}
Dengan menggunakan fungsi Heaviside, $f(t)$ dapat
ditulis sebagai
\begin{equation*}
f(t) = u(t - \pi) \, \sin t .
\end{equation*}
Demikian pula, fungsi tangga yang bernilai $1$ pada interval $[1,2)$ dan $0$
di tempat lain ditulis sebagai
\begin{equation*}
u(t - 1) - u(t-2) .
\end{equation*}
Dengan fungsi Heaviside, kita dapat menyatakan fungsi yang didefinisikan sepotong-sepotong.
Jika $f(t) = t$ ketika $t$ berada dalam $[0,1]$,
$f(t) = -t+2$
ketika $t$ berada dalam $[1,2]$, dan $f(t) = 0$ selain itu, maka kita menulis
\begin{equation*}
f(t) = t \, \bigl( u(t) - u(t-1) \bigr) + 
(-t+2) \, \bigl( u(t - 1) - u(t-2) \bigr) .
\end{equation*}

Bagaimana fungsi Heaviside berinteraksi dengan transformasi
Laplace?  Kita telah melihat bahwa (untuk $a \geq 0$)
\begin{equation*}
\mathcal{L} \bigl\{ u(t-a) \bigr\} = \frac{e^{-as}}{s} .
\end{equation*}
Perhitungan ini dapat digeneralisasi menjadi suatu \emph{\myindex{sifat pergeseran}}
atau \emph{\myindex{sifat pergeseran kedua}}:
\begin{equation} \label{ltd:sseq}
\mybxbg{~~
\mathcal{L} \bigl\{ f(t-a) \, u(t-a) \bigr\}
= e^{-as} \mathcal{L} \bigl\{ f(t) \bigr\}
\qquad (a \geq 0).
~~}
\end{equation}
Sebagai contoh,
\begin{equation*}
\mathcal{L} \bigl\{ t \, u(t-1) \bigr\}
=
\mathcal{L} \bigl\{ \bigl( (t-1) + 1\bigr) \, u(t-1) \bigr\}
= e^{-s} \mathcal{L} \bigl\{ t+1 \bigr\}
= e^{-s} \left(\frac{1}{s^2} + \frac{1}{s} \right)
.
\end{equation*}

\begin{example} \label{lt:rocketex}
Perhatikan
sistem massa-pegas
\begin{equation*}
x''(t) + x(t) = f(t) , \quad x(0) = 0, \quad x'(0) = 0,
\end{equation*}
dengan $f(t) = 1$ jika $1 \leq t < 5$ dan nol selain itu.
Fungsi $f(t)$ tidak periodik dan didefinisikan sepotong-sepotong; hal itu tidak menjadi masalah bagi Laplace.
Bayangkan sebuah
roket yang dipasang pada massa dinyalakan selama 4 detik mulai dari
$t=1$.  Atau, bayangkan suatu rangkaian RLC yang tegangannya dinaikkan
dengan laju konstan selama 4 detik mulai dari $t=1$, lalu dipertahankan tetap 
kembali
mulai dari $t=5$ (ingat bahwa $f(t)$ merepresentasikan turunan
tegangan dalam rangkaian RLC).

\pagebreak[2]
Kita menulis $f(t) = u(t-1) - u(t-5)$.  Kita mentransformasikan persamaan dan menyubstitusikan
syarat awal seperti sebelumnya
\begin{equation*}
s^2 X(s) + X(s) = \frac{e^{-s}}{s} - \frac{e^{-5s}}{s} .
\end{equation*}
Kita selesaikan untuk $X(s)$,
\begin{equation*}
X(s) = \frac{e^{-s}}{s(s^2+1)} - \frac{e^{-5s}}{s(s^2+1)} .
\end{equation*}
Kita serahkan sebagai latihan kepada pembaca untuk menunjukkan bahwa
\begin{equation*}
{\mathcal{L}}^{-1} \left\{ \frac{1}{s(s^2+1)} \right\}
= 1 - \cos t .
\end{equation*}
Dengan kata lain,
$\mathcal{L} \{ 1 - \cos t  \} = 
\frac{1}{s(s^2+1)}$.
Dengan menggunakan sifat pergeseran kedua \eqref{ltd:sseq}, kita memperoleh
\begin{equation*}
{\mathcal{L}}^{-1} \left\{ \frac{e^{-s}}{s(s^2+1)} \right\}
=
{\mathcal{L}}^{-1} \bigl\{
e^{-s}
\mathcal{L} \{ 1 - \cos t \}
\bigr\}
=
\bigl( 1 - \cos (t-1) \bigr) \, u(t-1) .
\end{equation*}
Demikian pula,
\begin{equation*}
{\mathcal{L}}^{-1} \left\{ \frac{e^{-5s}}{s(s^2+1)} \right\}
=
{\mathcal{L}}^{-1} \bigl\{
e^{-5s}
\mathcal{L} \{ 1 - \cos t \}
\bigr\}
=
\bigl( 1 - \cos (t-5) \bigr) \, u(t-5) .
\end{equation*}
Jadi, solusinya adalah
\begin{equation*}
x(t) = 
\bigl( 1 - \cos (t-1) \bigr) \, u(t-1) -
\bigl( 1 - \cos (t-5) \bigr) \, u(t-5) .
\end{equation*}
Grafik solusi ini diberikan dalam \figurevref{lt:heavisideexfig}.

\begin{myfig}
\capstart
\diffyincludegraphics{width=3in}{width=4.5in}{lt-heavisideex}{%
Grafik suatu fungsi untuk t dari 0 hingga 20. Fungsi bermula pada nol
dan tetap nol hingga t sama dengan 1, ketika grafik melengkung ke atas dan mengikuti kurva
sinusoidal yang sedikit terdeformasi sebelum menetap menjadi kurva sinusoidal biasa
setelah t sama dengan 5.}
\caption{Grafik $x(t)$.\label{lt:heavisideexfig}}
\end{myfig}
\end{example}

\subsection{Fungsi transfer}

Transformasi Laplace menghasilkan konsep berikut yang berguna untuk mempelajari
perilaku keadaan tunak suatu sistem linear.  Perhatikan persamaan
berbentuk
\begin{equation*}
L x = f(t) ,
\end{equation*}
dengan $L$ operator diferensial linear berkoefisien konstan.
Kemudian $f(t)$ biasanya dipandang sebagai masukan sistem dan $x(t)$
dipandang sebagai keluaran sistem.  Sebagai contoh, untuk sistem massa-pegas,
masukannya adalah fungsi pemaksa dan keluarannya adalah perilaku
massa.  Kita ingin memiliki cara yang mudah untuk mempelajari perilaku
sistem bagi berbagai masukan.

Andaikan bahwa
semua syarat awal bernilai nol.
Kita mengambil transformasi Laplace dari persamaan untuk memperoleh, bagi suatu $A(s)$,
\begin{equation*}
A(s) X(s) = F(s) .
\end{equation*}
Dengan menyelesaikan rasio $\nicefrac{X(s)}{F(s)}$, diperoleh apa yang disebut
\emph{\myindex{fungsi transfer}}
$H(s) = \nicefrac{1}{A(s)}$,
yakni,
\begin{equation*}
H(s) = \frac{X(s)}{F(s)} .
\end{equation*}
Dengan kata lain, $X(s) = H(s) F(s)$.  Kita memperoleh ketergantungan aljabar
keluaran sistem pada masukannya.  Sekarang kita dapat dengan mudah mempelajari
perilaku keadaan tunak sistem untuk berbagai masukan hanya dengan
mengalikan dengan fungsi transfer.
Selain itu, $H(s)$ dapat dihitung tanpa mengetahui secara tepat
persamaannya, yaitu dengan mengamati keluaran $X(s)$ untuk masukan tertentu $F(s)$.
Setelah $H(s)$ diketahui, Anda dapat mencari keluaran untuk sembarang masukan.

\begin{example}
Diberikan $x'' + \omega_0^2 x = f(t)$ (andaikan syarat awal bernilai nol),
mari kita cari fungsi transfernya.

Pertama, kita mengambil transformasi Laplace dari persamaan,
\begin{equation*}
s^2 X(s) + \omega_0^2 X(s) = F(s) .
\end{equation*}
Sekarang kita selesaikan fungsi transfer $\nicefrac{X(s)}{F(s)}$,
\begin{equation*}
H(s) = \frac{X(s)}{F(s)} = \frac{1}{s^2 + \omega_0^2} .
\end{equation*}

Mari kita lihat cara menggunakan fungsi transfer.  Andaikan kita memiliki masukan konstan
$f(t) = 1$.  Maka $F(s) = \nicefrac{1}{s}$, dan
\begin{equation*}
X(s) = H(s) F(s) = \frac{1}{s^2+\omega_0^2} \, \frac{1}{s} .
\end{equation*}
Dengan mengambil transformasi Laplace invers dari $X(s)$, kita memperoleh
\begin{equation*}
x(t) = \frac{1-\cos(\omega_0 t)}{\omega_0^2} .
\end{equation*}
Demikian pula, untuk sembarang masukan lain $F(s)$, keluarannya adalah
$X(s) = H(s) F(s) = \frac{1}{s^2+\omega_0^2} \, F(s)$.
\end{example}

\subsection{Transformasi integral}

Salah satu keunggulan transformasi Laplace adalah kemampuannya menangani
persamaan integral.  Ini adalah persamaan yang memuat integral, bukan
turunan fungsi.  Sifat dasar, yang dapat dibuktikan
dengan menerapkan definisi dan melakukan integrasi parsial, adalah 
\begin{equation*}
\mybxbg{~~
\mathcal{L} \left\{
\int_0^t f(\tau) \, d\tau
\right\} = \frac{1}{s} \, F(s) .
~~}
\end{equation*}
Kadang-kadang berguna (misalnya,\ untuk menghitung transformasi invers) menulis
hubungan ini sebagai
\begin{equation*}
\int_0^t f(\tau) \, d\tau
=
{\mathcal{L}}^{-1} \left\{
\frac{1}{s} \, F(s) \right\} .
\end{equation*}

\begin{example}
Untuk menghitung ${\mathcal{L}}^{-1} \left\{\frac{1}{s(s^2+1)}\right\}$, kita dapat
menerapkan aturan pengintegralan ini:
\begin{equation*}
{\mathcal{L}}^{-1} \left\{
\frac{1}{s} \, \frac{1}{s^2+1} \right\} 
=
\int_0^t 
{\mathcal{L}}^{-1} \left\{
\frac{1}{s^2+1} \right\} \, d\tau
=
\int_0^t 
\sin \tau \, d\tau
=
1 - \cos t .
\end{equation*}
\end{example}


\begin{example}
Persamaan yang memuat integral dari fungsi yang tidak diketahui
disebut \emph{\myindex{persamaan integral}}.
Perhatikan
\begin{equation*}
x(t) - t = \int_0^t x(\tau) \, d\tau ,
\end{equation*}
dengan kita ingin menyelesaikan untuk $x(t)$.
Kita menerapkan transformasi Laplace untuk memperoleh
\begin{equation*}
X(s) - \frac{1}{s^2} = \frac{1}{s} X(s) ,
\end{equation*}
dengan $X(s) = \mathcal{L} \bigl\{ x(t) \bigr\}$.  Jadi,
\begin{equation*}
X(s) = \frac{1}{s(s-1)} = \frac{1}{s-1} - \frac{1}{s} .
\end{equation*}
Transformasi Laplace invers memberikan
\begin{equation*}
x(t) = e^t - 1 .
\end{equation*}
%More complicated integral
%equations can also be solved using convolution, about which we will
%learn in the next section.
\end{example}

\subsection{Fungsi periodik}

Pembaca mungkin bertanya: Bagaimana dengan fungsi periodik sebagai masukan $f(t)$?
Artinya, suatu fungsi $f(t)$ dengan $f(t) = f(t+P)$ untuk suatu konstanta $P$ (periodenya).
Baiklah, mari kita hitung $F(s)$:
\begin{equation*}
\begin{split}
F(s) & = \int_0^\infty e^{-st} f(t) \, dt
%mbxlatex \\
%mbxlatex &
= \int_0^P e^{-st} f(t) \, dt
+ \int_P^\infty e^{-st} f(t) \, dt
\\
& = \int_0^P e^{-st} f(t) \, dt
+ \int_0^\infty e^{-s(t+P)} f(t+P) \, dt
%mbxlatex \\
%mbxlatex &
= \int_0^P e^{-st} f(t) \, dt
+ e^{-Ps} \int_0^\infty e^{-st} f(t) \, dt
\\
& = \int_0^P e^{-st} f(t) \, dt
+ e^{-Ps} F(s) .
\end{split}
\end{equation*}
Dengan menyelesaikan untuk $F(s)$, kita memperoleh
\begin{equation*}
F(s)
=
\frac{1}{1-e^{-Ps}}
\int_0^P e^{-st} f(t) \, dt .
\end{equation*}
Seperti sebelumnya, menghitung inversnya akan lebih rumit dan mungkin
melibatkan pencarian dalam tabel.  Kita tidak perlu membahas perhitungan invers
di sini.

\begin{example}
Andaikan $f(t)$ adalah suatu versi \emph{\myindex{gelombang gigi gergaji}}, yakni,
misalkan
$f(t) = t$ untuk $0 \leq t < 1$
dan gunakan $f(t)=f(t+1)$ untuk memperluasnya secara periodik.  Jadi,
$f(t) = t-1$ untuk $1 \leq t < 2$, 
$f(t) = t-2$ untuk $2 \leq t < 3$, dan seterusnya.  Kemudian $P=1$ dan
perhitungan singkat dengan integrasi parsial memberikan
\begin{equation*}
F(s) =
\frac{1}{1-e^{-s}}
\int_0^1 e^{-st} t \, dt 
=
\frac{1}{1-e^{-s}}
\left(
\frac{-e^{-s}}{s}-\frac{e^{-s}}{s^2} + \frac{1}{s^2}
\right)
=
\frac{-e^{-s}}{(1-e^{-s})s} + \frac{1}{s^2} .
\end{equation*}
\end{example}

\subsection{Latihan}

\begin{exercise}
Dengan menggunakan fungsi Heaviside, tuliskan fungsi sepotong-sepotong
yang bernilai 0 untuk $t < 0$, $t^2$ untuk $t$ dalam $[0,1]$, dan $t$ untuk $t > 1$.
\end{exercise}

\begin{exercise}
Dengan menggunakan transformasi Laplace, selesaikan
\begin{equation*}
m x'' + c x' + k x = 0 , \quad x(0) = a, \quad x'(0) = b ,
\end{equation*}
dengan $m > 0$, $c > 0$, $k > 0$, dan
$c^2 - 4km > 0$ (sistem teredam berlebih).
\end{exercise}

\begin{exercise}
Dengan menggunakan transformasi Laplace, selesaikan
\begin{equation*}
m x'' + c x' + k x = 0 , \quad x(0) = a, \quad x'(0) = b ,
\end{equation*}
dengan $m > 0$, $c > 0$, $k > 0$, dan
$c^2 - 4km < 0$ (sistem teredam kurang).
\end{exercise}

\begin{exercise}
Dengan menggunakan transformasi Laplace, selesaikan
\begin{equation*}
m x'' + c x' + k x = 0 , \quad x(0) = a, \quad x'(0) = b ,
\end{equation*}
dengan $m > 0$, $c > 0$, $k > 0$, dan
$c^2 = 4km$ (sistem teredam kritis).
\end{exercise}

\begin{exercise}
Selesaikan $x'' + x = u(t-1)$ dengan syarat awal $x(0) = 0$ dan $x'(0) = 0$.
\end{exercise}

\begin{exercise}
Tunjukkan sifat \myquote{pendiferensialan transformasi}.  Andaikan
$\mathcal{L} \bigl\{ f(t) \bigr\} = F(s)$, lalu tunjukkan
\begin{equation*}
\mathcal{L} \bigl\{ -t f(t) \bigr\} = F'(s) .
\end{equation*}
Petunjuk: Diferensialkan di bawah tanda integral.
\end{exercise}

\begin{exercise}
Selesaikan $x''' + x = t^3 u(t-1)$ dengan syarat awal $x(0) = 1$ dan $x'(0) =
0$, $x''(0) = 0$.
\end{exercise}

\begin{exercise}
Dengan menggunakan definisi transformasi Laplace,
justifikasikan sifat pergeseran kedua: 
$\mathcal{L} \bigl\{ f(t-a) \, u(t-a) \bigr\}
= e^{-as} \mathcal{L} \bigl\{ f(t) \bigr\}$.
\end{exercise}

\begin{exercise}
Perhatikan sistem massa-pegas dengan roket dari
\exampleref{lt:rocketex}.  Kita melihat bahwa solusi terus berosilasi
setelah roket berhenti menyala.  Amplitudo osilasi bergantung
pada lamanya roket dinyalakan (selama 4 detik dalam contoh).
\begin{tasks}
\task
Carilah rumus untuk amplitudo osilasi yang dihasilkan
dalam hubungannya dengan lama waktu roket dinyalakan.
\task
 Apakah terdapat
waktu taknol (jika ya, berapa?)
ketika roket menyala dan osilasi yang dihasilkan
memiliki amplitudo 0 (massa tidak bergerak)?
\end{tasks}
\end{exercise}

\begin{exercise}
Definisikan
\begin{equation*}
f(t) =
\begin{cases}
{(t-1)}^2 & \text{jika } \; 1 \leq t < 2, \\
3-t & \text{jika } \; 2 \leq t < 3, \\
0 & \text{selain itu} .
\end{cases}
\end{equation*}
\begin{tasks}
\task Buat sketsa grafik $f(t)$.
\task Tuliskan $f(t)$ menggunakan fungsi Heaviside.
\task Selesaikan $x''+x=f(t)$, $x(0)=0$, $x'(0) = 0$ menggunakan transformasi Laplace.
\end{tasks}
\end{exercise}

\begin{exercise}
Carilah fungsi transfer untuk 
$m x'' + c x' + kx = f(t)$,
$x(0)=0$, $x'(0)=0$.
\end{exercise}

\begin{exercise}
Andaikan $Lx = f(t)$, $x(0)=0$, $x'(0)=0$
untuk fungsi masukan $f(t) = t$ memberikan keluaran
$x(t) = 2e^{-t} + te^{-t} + (t-2)$.
\begin{tasks}
\task Carilah $F(s)$, $X(s)$, dan fungsi transfer $H(s)$.
\task Jika masukannya adalah $f(t) = \sin(t)$, carilah keluaran baru $x(t)$.
\end{tasks}
\end{exercise}

\begin{exercise}
Andaikan $f(t) = 1$ jika $0 \leq t < 1$ dan $f(t)=0$ jika $1 \leq t < 2$, lalu
perluas secara periodik untuk semua $t \geq 0$ sehingga $f(t)=f(t+2)$.  Hitung
transformasi Laplace $F(s)$.
\end{exercise}

\setcounter{exercise}{100}

\begin{exercise}
Dengan menggunakan fungsi Heaviside $u(t)$, tuliskan fungsi
\begin{equation*}
f(t) =
\begin{cases}
0 & \text{jika } \; \phantom{1 \leq {}} t < 1  , \\
t-1 & \text{jika } \; 1 \leq t < 2 , \\
1 & \text{jika } \; 2 \leq t .
\end{cases}
\end{equation*}
\end{exercise}
\exsol{%
$f(t) =
(t-1)\bigl(u(t-1) - u(t-2)\bigr) + 
u(t-2)$
}

\begin{exercise}
Selesaikan $x''-x = (t^2-1) u(t-1)$ dengan syarat awal $x(0)=1$, $x'(0) = 2$
menggunakan transformasi Laplace.
\end{exercise}
\exsol{%
$x(t) = (2e^{t-1}-t^2-1) u(t-1) -\frac{1}{2}e^{-t}+\frac{3}{2}e^t$
}
 
\begin{exercise}
Carilah fungsi transfer untuk 
$x' + x = f(t)$,
$x(0)=0$.
\end{exercise}
\exsol{%
$H(s) = \frac{1}{s+1}$
}

\begin{exercise}
Andaikan $Lx = f(t)$, $x(0)=0$, $x'(0)=0$
untuk fungsi masukan $f(t) = 4$ memberikan keluaran
$x(t) = 1 - \cos(2 t)$.  Carilah $F(s)$, $X(s)$, dan fungsi transfer
$H(s)$.
\end{exercise}
\exsol{%
$F(s) = \frac{4}{s}$,
$X(s) = \frac{1}{s} - \frac{s}{s^2+4}$,
$H(s) = \frac{1}{s^2+4}$.
}

%%%%%%%%%%%%%%%%%%%%%%%%%%%%%%%%%%%%%%%%%%%%%%%%%%%%%%%%%%%%%%%%%%%%%%%%%%%%%%

\sectionnewpage
\section{Konvolusi}
\label{convolution:section}

%mbxINTROSUBSECTION

\sectionnotes{1 atau 1.5 kuliah\EPref{, \S7.2 dalam \cite{EP}}\BDref{,
\S6.6 dalam \cite{BD}}}

\subsection{Konvolusi}

Transformasi Laplace dari suatu hasil kali bukanlah hasil kali
transformasi-transformasinya.  Namun, harapan belum hilang.  Kita hanya perlu menggunakan
jenis \myquote{hasil kali} yang berbeda.
Ambil
dua fungsi $f(t)$ dan $g(t)$ yang didefinisikan untuk $t \geq 0$,
dan definisikan \emph{\myindex{konvolusi}}%
\footnote{%
Bagi yang pernah melihat konvolusi sebelumnya, Anda mungkin pernah
melihatnya didefinisikan sebagai
$(f * g)(t) =
\int_{-\infty}^\infty f(\tau) g(t-\tau) \, d\tau$.  Definisi ini
bersepakat dengan \eqref{ltc:convdef} jika Anda mendefinisikan $f(t)$ dan $g(t)$
bernilai nol untuk $t < 0$.
Ketika membahas transformasi Laplace, definisi yang kita berikan sudah
memadai.  Namun, konvolusi muncul dalam banyak penerapan lain,
yang mungkin mengharuskan Anda menggunakan definisi lebih umum dengan batas takhingga.
}
dari $f(t)$ dan $g(t)$ sebagai
\begin{equation} \label{ltc:convdef}
\mybxbg{~~
(f * g)(t) \overset{\text{def}}{=}
\int_0^t f(\tau) g(t-\tau) \, d\tau .
~~}
\end{equation}
Seperti yang dapat dilihat, konvolusi dua fungsi dari $t$ adalah fungsi lain dari $t$.


\begin{example}
Ambil $f(t) = e^t$ dan $g(t) = t$ untuk $t \geq 0$.  Maka 
\begin{equation*}
(f*g)(t)
=
\int_0^t e^\tau (t-\tau) \, d\tau
=
e^t - t - 1 .
\end{equation*}
Untuk menyelesaikan integral tersebut, kita
melakukan satu kali integrasi parsial.
\end{example}

\begin{example} \label{ltc:convsincosex}
Ambil $f(t) = \sin (\omega t)$ dan $g(t) = \cos (\omega t)$ untuk $t \geq 0$.
Maka 
\begin{equation*}
(f*g)(t)
=
\int_0^t  \sin ( \omega \tau ) \,
\cos \bigl( \omega (t-\tau) \bigr) \, d\tau .
\end{equation*}
Terapkan identitas
\begin{equation*}
\cos (\theta) \sin (\psi) =
\frac{1}{2} \, \bigl( \sin (\theta + \psi) - \sin (\theta - \psi) \bigr) ,
\end{equation*}
untuk memperoleh
\begin{equation*}
\begin{split}
(f*g)(t)
& =
\int_0^t
\frac{1}{2} \, \bigl( \sin (\omega t) - \sin (\omega t - 2 \omega \tau
) \bigr) \, d\tau
\\
& =
\left[ \frac{1}{2} \, \tau  \sin (\omega t) - \frac{1}{4\omega} \, \cos
(\omega t - 2 \omega \tau) \right]_{\tau=0}^t
\\
& = \frac{1}{2} \, t \sin (\omega t) .
\end{split}
\end{equation*}
Rumus tersebut hanya berlaku untuk $t \geq 0$.  Fungsi $f$, $g$,
dan $f*g$ tidak terdefinisi untuk $t < 0$.
\end{example}

Konvolusi memiliki banyak sifat yang membuatnya berperilaku seperti hasil kali.
Misalkan $c$ adalah konstanta dan $f$, $g$, serta $h$ adalah fungsi.
Memverifikasikan bahwa
\begin{align*}
& f * g = g * f , \\
& (c f) * g = f * (c g) = c (f*g) , \\
& (f+g) * h = f * h + g * h , \\
& ( f * g ) * h = f * ( g * h ) .
\end{align*}
merupakan latihan kalkulus. Sifat pertama berarti bahwa $(f * g)(t) =
\int_0^t f(\tau) g(t-\tau) \, d\tau = 
\int_0^t f(t-\tau) g(\tau) \, d\tau$.
Sifat yang paling menarik bagi kita adalah teorema berikut.

\begin{theorem}
Jika $f(t)$ dan $g(t)$ berorde eksponensial, maka
$(f*g)(t)$ juga demikian dan
\begin{equation*}
\mybxbg{~~
\mathcal{L} \bigl\{ (f*g)(t) \bigr\}
=
\mathcal{L} \left\{ \int_0^t f(\tau) g(t-\tau) \, d\tau \right\}
=
\mathcal{L} \bigl\{ f(t) \bigr\} \mathcal{L} \bigl\{ g(t) \bigr\} .
~~}
\end{equation*}
\end{theorem}

Dengan kata lain, transformasi Laplace dari suatu konvolusi adalah hasil kali
transformasi-transformasi Laplace:
${\mathcal{L}} \bigl\{ (f*g)(t) \bigr\} = F(s)G(s)$,
atau sebaliknya,
${\mathcal{L}}^{-1} \bigl\{ F(s) G(s) \bigr\} = (f*g)(t)$.

\begin{example}
Andaikan kita ingin mencari transformasi Laplace invers dari
\begin{equation*}
\frac{1}{s^2(s+1)} = 
\frac{1}{s^2}
\,
\frac{1}{s+1}
.
\end{equation*}
Kita mengenali dua entri dari \tableref{ltd:table}:
\begin{equation*}
\mathcal{L}^{-1} 
\left\{
\frac{1}{s^2} \right\} 
= t
\qquad \text{dan} \qquad
\mathcal{L}^{-1} 
\left\{
\frac{1}{s+1} \right\}
= e^{-t}
.
\end{equation*}
Oleh karena itu, kita mengonvolusikan $t$ dan $e^{-t}$,
\begin{equation*}
\mathcal{L}^{-1}
\left\{
\frac{1}{s^2}
\,
\frac{1}{s+1}
\right\}
=
\int_0^t
\tau e^{-(t-\tau)} \,d\tau
=
e^{-t}+t-1 .
\end{equation*}
Perhitungan integral tersebut melibatkan integrasi parsial.
\end{example}

\subsection{Menyelesaikan ODE}

Contoh berikut menunjukkan kekuatan penuh konvolusi dan
transformasi Laplace.  Kita dapat memberikan solusi bagi
masalah osilasi paksa untuk sembarang fungsi pemaksa sebagai integral
tentu.

\begin{example} \label{example:undampedforcedbylaplacearbitrhs}
Carilah solusi bagi
\begin{equation*}
x'' + \omega_0^2 x = f(t) , \quad x(0) = 0, \quad x'(0) = 0 ,
\avoidbreak
\end{equation*}
untuk sembarang fungsi $f(t)$.

Pertama, kita menerapkan transformasi Laplace pada persamaan.  Nyatakan
transformasi dari $x(t)$ dengan $X(s)$ dan transformasi dari $f(t)$ dengan
$F(s)$ seperti biasa.  Kita memperoleh
\begin{equation*}
s^2 X(s) + \omega_0^2 X(s) = F(s) ,
\end{equation*}
atau dengan kata lain,
\begin{equation*}
X(s) = \frac{1}{s^2+ \omega_0^2} F(s).
\end{equation*}
Ingat bahwa $H(s) = \frac{1}{s^2+ \omega_0^2}$ adalah fungsi transfer.
Kita mengetahui
\begin{equation*}
{\mathcal{L}}^{-1} \left\{
\frac{1}{s^2+ \omega_0^2}
\right\} = 
\frac{\sin (\omega_0 t)}{\omega_0} .
\end{equation*}
Oleh karena itu,
\begin{equation*}
x(t) = 
\int_0^t
\frac{\sin (\omega_0 \tau)}{\omega_0}
f(t-\tau) \, d\tau ,
\end{equation*}
atau, jika kita membalik urutannya,
\begin{equation*}
x(t) = 
\int_0^t
f(\tau) 
\frac{\sin \bigl( \omega_0 (t-\tau) \bigr)}{\omega_0} \, d\tau .
\end{equation*}
\end{example}

Perhatikan satu lagi ciri dari contoh di atas.
Sekarang kita dapat melihat cara transformasi Laplace
menangani \myindex{resonansi}.  Andaikan $f(t) =
\cos (\omega_0 t)$.  Maka
\begin{equation*}
x(t) = 
\int_0^t
\frac{\sin (\omega_0 \tau)}{\omega_0} \,
\cos \bigl( \omega_0 (t-\tau) \bigr) \, d\tau
=
\frac{1}{\omega_0}
\int_0^t
\sin ( \omega_0 \tau ) \,
\cos \bigl(\omega_0 (t-\tau) \bigr) \, d\tau .
\end{equation*}
Kita telah menghitung konvolusi sinus dan kosinus dalam
\exampleref{ltc:convsincosex}.  Jadi,
\begin{equation*}
x(t) =
\left(
\frac{1}{\omega_0}
\right) \,
\left(
\frac{1}{2} \,
t
\sin ( \omega_0 t )
\right)
=
\frac{1}{2 \omega_0} \,
t
\sin ( \omega_0 t ).
\end{equation*}
Perhatikan $t$ di depan sinus.  Karena itu, solusi bertumbuh tanpa
batas ketika $t$ membesar, yang berarti terjadi resonansi.

Gagasan umumnya di sini adalah bahwa jika $H(s)$ merupakan fungsi transfer, maka
$X(s)=H(s)F(s)$.
Jika kita mencari $h(t) = {\mathcal{L}}^{-1}\bigl\{ H(s) \bigr\}$, maka
\begin{equation*}
x(t)
= {\mathcal{L}}^{-1}\bigl\{ X(s) \bigr\}
= {\mathcal{L}}^{-1}\bigl\{ F(s)H(s) \bigr\}
= (f * h)(t)
= \int_0^t f(\tau) h(t-\tau) \, d\tau .
\end{equation*}
Jadi,
kita dapat menyelesaikan setiap persamaan berkoefisien konstan dengan sembarang fungsi
pemaksa $f(t)$ sebagai integral tentu menggunakan konvolusi.
Integral tentu, alih-alih solusi berbentuk tertutup, biasanya sudah memadai
untuk kebanyakan keperluan praktis.  Menghitung integral tentu secara numerik
tidaklah sulit.

\subsection{Persamaan integral Volterra}

Suatu persamaan integral\index{persamaan integral}
yang melibatkan konvolusi
dan muncul, misalnya, ketika mempelajari sistem dengan efek ingatan
adalah \emph{\myindex{persamaan integral Volterra}}%
\footnote{Dinamai menurut matematikawan Italia
\href{https://en.wikipedia.org/wiki/Vito_Volterra}{Vito Volterra}
(1860--1940).}
\begin{equation*}
x(t) = f(t) + \int_0^t g(t-\tau) x(\tau) \, d\tau .
\end{equation*}
Di sini $f(t)$ dan $g(t)$ adalah fungsi yang diketahui dan $x(t)$ adalah fungsi takdiketahui yang
ingin kita cari.
Kita memiliki $x(t) = f(t) + (g * x)(t)$, jadi pertanyaannya adalah:
Jika konvolusi
menyerupai hasil kali, dapatkah kita juga \myquote{membagi} untuk menyelesaikan persamaan ini?
Kita menerapkan transformasi Laplace pada persamaan untuk memperoleh 
\begin{equation*}
X(s) = F(s) + G(s) X(s) ,
\end{equation*}
dengan $X(s)$, $F(s)$, dan $G(s)$ masing-masing transformasi Laplace dari $x(t)$, $f(t)$, dan
$g(t)$.  Kita memperoleh
\begin{equation*}
X(s) = \frac{F(s)}{1-G(s)} .
\end{equation*}
Untuk mencari $x(t)$, sekarang kita perlu mencari 
transformasi Laplace invers dari $X(s)$.

\begin{example}
Selesaikan
\begin{equation*}
x(t) =  e^{-t} + \int_0^t \sinh(t-\tau) x(\tau) \, d\tau .
\end{equation*}

Kita menerapkan transformasi Laplace untuk memperoleh
\begin{equation*}
X(s) = \frac{1}{s+1} + \frac{1}{s^2-1} X(s) ,
\end{equation*}
atau
\begin{equation*}
X(s) = \frac{\frac{1}{s+1}}{1- \frac{1}{s^2-1}}
=
\frac{s-1}{s^2 - 2}
=
\frac{s}{s^2 - 2}
-
\frac{1}{s^2 - 2} .
\end{equation*}
Tidak sulit menerapkan \tablevref{lt:table} untuk memperoleh
\begin{equation*}
x(t) = \cosh \bigl( \sqrt{2} \, t \bigr) -
\frac{1}{\sqrt{2}} \sinh \bigl( \sqrt{2}\, t \bigr).
\end{equation*}
\end{example}

\subsection{Latihan}

\begin{exercise}
Misalkan $f(t) = t^2$ untuk $t \geq 0$, dan $g(t) = u(t-1)$.  Hitung
$f * g$.
\end{exercise}

\begin{exercise}
Misalkan $f(t) = t$ untuk $t \geq 0$, dan $g(t) = \sin t $ untuk $t \geq 0$.  Hitung
$f * g$.
\end{exercise}

\begin{exercise}
Carilah solusi bagi
\begin{equation*}
m x'' + c x' + k x = f(t) , \quad x(0) = 0, \quad x'(0) = 0 ,
\end{equation*}
untuk sembarang fungsi $f(t)$, dengan $m > 0$, $c > 0$, $k > 0$,
dan $c^2 - 4km > 0$ (sistem teredam berlebih).
Tuliskan solusi sebagai integral tentu.
\end{exercise}

\begin{exercise}
Carilah solusi bagi
\begin{equation*}
m x'' + c x' + k x = f(t) , \quad x(0) = 0, \quad x'(0) = 0 ,
\end{equation*}
untuk sembarang fungsi $f(t)$, dengan $m > 0$, $c > 0$, $k > 0$,
dan $c^2 - 4km < 0$ (sistem teredam kurang).
Tuliskan solusi sebagai integral tentu.
\end{exercise}

\begin{exercise}
Carilah solusi bagi
\begin{equation*}
m x'' + c x' + k x = f(t) , \quad x(0) = 0, \quad x'(0) = 0 ,
\end{equation*}
untuk sembarang fungsi $f(t)$, dengan $m > 0$, $c > 0$, $k > 0$,
dan $c^2 = 4km$ (sistem teredam kritis).
Tuliskan solusi sebagai integral tentu.
\end{exercise}

\begin{exercise}
Selesaikan
\begin{equation*}
x(t) =  e^{-t} + \int_0^t \cos(t-\tau) x(\tau) \, d\tau .
\end{equation*}
\end{exercise}

\begin{exercise}
Selesaikan
\begin{equation*}
x(t) =  \cos t + \int_0^t \cos(t-\tau) x(\tau) \, d\tau .
\end{equation*}
\end{exercise}

\begin{exercise}
Hitung ${\mathcal{L}}^{-1} \left\{ \frac{s}{{(s^2+4)}^2} \right\}$ menggunakan
konvolusi.
\end{exercise}

\begin{exercise}
Tuliskan solusi bagi
$x''-2x=e^{-t^2}$, $x(0)=0$, $x'(0)=0$ sebagai
integral tentu.  Petunjuk: Jangan mencoba menghitung
transformasi Laplace dari $e^{-t^2}$.
\end{exercise}

\setcounter{exercise}{100}

\begin{exercise}
Misalkan $f(t) = \cos t$ untuk $t \geq 0$, dan $g(t) = e^{-t}$.  Hitung
$f * g$.
\end{exercise}
\exsol{%
$\frac{1}{2}(\cos t + \sin t - e^{-t})$
}

\begin{exercise}
Hitung ${\mathcal{L}}^{-1} \left\{ \frac{5}{s^4+s^2} \right\}$ menggunakan
konvolusi.
\end{exercise}
\exsol{%
$5t-5\sin t$
}

\begin{exercise}
Selesaikan $x''+x = \sin t$, $x(0) = 0$, $x'(0)=0$ menggunakan konvolusi.
\end{exercise}
\exsol{%
$\frac{1}{2}(\sin t - t \cos t)$
}

\begin{exercise}
Selesaikan $x'''+x' = f(t)$, $x(0) = 0$, $x'(0)=0$, $x''(0)=0$ menggunakan konvolusi.
Tuliskan hasilnya sebagai integral tentu.
\end{exercise}
\exsol{%
$\int_0^t f(\tau) \bigl( 1 - \cos (t-\tau)\bigr)\, d\tau$
}

%%%%%%%%%%%%%%%%%%%%%%%%%%%%%%%%%%%%%%%%%%%%%%%%%%%%%%%%%%%%%%%%%%%%%%%%%%%%%%

\sectionnewpage
\section{Delta Dirac dan respons impuls}
\label{diracdelta:section}

%mbxINTROSUBSECTION

\sectionnotes{1 atau 1.5 kuliah\EPref{, \S7.6 dalam \cite{EP}}\BDref{,
\S6.5 dalam \cite{BD}}}

\subsection{Pulsa persegi panjang}

Sering kali kita mempelajari suatu sistem fisik dengan memberikan pulsa singkat 
dan kemudian melihat apa yang dilakukan sistem tersebut.
Perilaku yang dihasilkan
sering disebut \emph{\myindex{respons impuls}},
dan memahaminya memberi tahu kita cara sistem merespons sembarang masukan.
Mari kita lihat apa yang dimaksud dengan pulsa.
Jenis pulsa paling sederhana adalah
\emph{\myindex{pulsa persegi panjang}}\index{pulsa}
yang didefinisikan oleh
\begin{equation*}
\varphi(t) = 
\begin{cases}
0 & \text{jika } \; \phantom{a \leq {}} t < a , \\
M & \text{jika } \; a \leq t < b , \\
0 & \text{jika } \; b \leq t .
\end{cases}
\end{equation*}
Lihat \figurevref{lt:sqpulse} untuk grafiknya.

%14 is number of lines, must be adjusted
\begin{mywrapfig}[14]{3.25in}
\capstart
\diffyincludegraphics{width=3in}{width=4.5in}{lt-sqpulse}{%
Grafik pulsa persegi panjang yang ditampilkan dari t sama dengan 0 hingga t sama dengan 3.
Fungsi bernilai nol hingga t sama dengan 0.5, kemudian bernilai 2
antara t sama dengan 0.5 dan t sama dengan 1.0, lalu kembali bernilai nol sejak
t sama dengan 1.0 dan seterusnya.}
\caption{Contoh pulsa persegi panjang dengan $a=0.5$, $b=1$, dan $M = 2$.\label{lt:sqpulse}}
\end{mywrapfig}

Perhatikan bahwa
\begin{equation*}
\varphi(t) = M \bigl( u(t-a) - u(t-b) \bigr) ,
\end{equation*}
dengan $u(t)$ fungsi tangga satuan.
Transformasi Laplace dari pulsa persegi panjang adalah
\begin{equation*}
\begin{split}
{\mathcal{L}} \bigl\{ \varphi(t) \bigr\}
& =
{\mathcal{L}} \bigl\{ M \bigl( u(t-a) - u(t-b) \bigr)  \bigr\}
\\
& =
M
\frac{e^{-as} - e^{-bs}}{s} .
\end{split}
\end{equation*}

Demi kesederhanaan, tetapkan $a=0$.
Kita juga
menetapkan $M = \nicefrac{1}{b}$ sehingga
\begin{equation*}
\int_0^\infty \varphi(t) \,dt = 1 .
\end{equation*}
Artinya, kita memiliki pulsa dengan \myquote{massa satuan.}
Dengan cara ini, untuk setiap $b$, pulsa tersebut memiliki \myquote{daya dorong} yang sama.
Untuk pulsa semacam itu,
\begin{equation*}
{\mathcal{L}} \bigl\{ \varphi(t) \bigr\}
=
{\mathcal{L}} \left\{ \frac{u(t) - u(t-b)}{b}  \right\}
=
\frac{1 - e^{-bs}}{bs} .
\end{equation*}
Kita menginginkan $b$ sangat kecil; kita ingin agar
pulsanya sangat singkat dan sangat tinggi.  Dengan membiarkan $b$ menuju nol, kita sampai
pada konsep fungsi delta Dirac.  Limit
dari $\frac{1-e^{-bs}}{bs}$
ketika $b \to 0$ adalah $1$, jadi kita mencari suatu \myquote{fungsi} yang transformasi
Laplace-nya adalah $1$.

\subsection{Fungsi delta}

\emph{\myindex{Fungsi delta Dirac}}\index{fungsi delta}%
\footnote{Dinamai menurut fisikawan dan matematikawan Inggris
\href{https://en.wikipedia.org/wiki/Paul_Dirac}{Paul Adrien Maurice Dirac}
(1902--1984).}
tidak persis berupa fungsi; objek ini kadang-kadang disebut
\emph{\myindex{fungsi tergeneralisasi}}.  Kita
menghindari perincian yang tidak perlu dan cukup mengatakan bahwa ini adalah suatu objek
yang sebenarnya tidak bermakna kecuali jika kita mengintegralkannya.  Motivasinya adalah
bahwa kita menginginkan suatu \myquote{fungsi} $\delta(t)$
sedemikian sehingga 
untuk setiap fungsi kontinu $f(t)$,
\begin{equation*}
\mybxbg{~~
\int_{-\infty}^\infty \delta(t) f(t) \,dt = f(0) .
~~}
\end{equation*}
Rumus tersebut seharusnya berlaku jika kita mengintegralkan pada sembarang interval yang memuat 0,
bukan hanya $(-\infty,\infty)$.
Jadi, $\delta(t)$ adalah suatu \myquote{fungsi} 
dengan seluruh \myquote{massanya} berada di satu titik $t=0$.  Untuk sembarang
interval\footnote{%
Penting bahwa kita menganggap $c$ dan $d$ sebagai bagian dari
interval.  Integral ini dapat ditulis sebagai $\int_{c-}^{d+}$.}
$[c,d]$,
\begin{equation*}
\int_c^d \delta(t) \,dt = 
\begin{cases}
1 & \text{jika interval $[c,d]$ memuat 0, yakni\ } c \leq 0 \leq d, \\
0 & \text{selain itu.}
\end{cases}
\end{equation*}
Sayangnya, tidak ada fungsi semacam itu dalam arti klasik.  Secara
informal, Anda dapat membayangkan bahwa $\delta(t)$ bernilai nol untuk $t \neq 0$ dan entah bagaimana
takhingga di $t=0$.

Cara yang baik untuk memandang $\delta(t)$ adalah sebagai limit pulsa
dengan panjang yang semakin kecil,
dan masing-masing memiliki integral $1$.  Sebagai contoh, 
perhatikan pulsa persegi panjang
$\varphi(t)$ seperti di atas dengan $a=0$ dan
$M=\nicefrac{1}{b}$, yakni, $\varphi(t) = \frac{u(t) - u(t-b)}{b}$.
Hitung
\begin{equation*}
\int_{-\infty}^\infty \varphi(t) f(t) \,dt =
\int_{-\infty}^\infty \frac{u(t) - u(t-b)}{b} f(t) \,dt =
\frac{1}{b} \int_{0}^b f(t) \,dt .
\end{equation*}
Jika $f(t)$ kontinu di $t=0$, maka
untuk $b$ yang sangat kecil, fungsi $f(t)$ kira-kira sama dengan $f(0)$ pada
interval $[0,b]$.  Kita aproksimasikan integral tersebut:
\begin{equation*}
\frac{1}{b} \int_{0}^b f(t) \,dt \approx
\frac{1}{b} \int_{0}^b f(0) \,dt = f(0) .
\end{equation*}
Jadi,
\begin{equation*}
\lim_{b\to 0}
\int_{-\infty}^\infty \varphi(t) f(t) \,dt =
\lim_{b\to 0}
\frac{1}{b} \int_{0}^b f(t) \,dt  = f(0) .
\end{equation*}

Karena itu, mari kita terima $\delta(t)$ sebagai objek yang dapat
diintegralkan.  Kita sering ingin menggeser $\delta$ ke titik lain, misalnya
$\delta(t-a)$.  Dalam hal ini,
\begin{equation*}
\int_{-\infty}^\infty \delta(t-a) f(t) \,dt = f(a) .
\end{equation*}
Perhatikan bahwa $\delta(a-t)$ adalah objek yang sama dengan $\delta(t-a)$.
Dengan demikian, konvolusi $\delta(t)$ dengan $f(t)$ kembali berupa $f(t)$,
\begin{equation*}
(f * \delta) (t) = 
\int_{0}^t f(\tau) \delta(t-\tau) \,d\tau
= f(t) .
\end{equation*}

Karena kita dapat mengintegralkan $\delta(t)$, kita menghitung transformasi Laplace-nya:
\begin{equation*}
\mybxbg{~~
{\mathcal{L}} \bigl\{ \delta(t-a) \bigr\}
=
\int_{0}^\infty e^{-st} \delta(t-a) \,dt = e^{-as} .
~~}
\end{equation*}
Secara khusus,
\begin{equation*}
{\mathcal{L}} \bigl\{ \delta(t) \bigr\} = 1 .
\end{equation*}

\begin{remark} \label{remark:deltaderivofjump}
Transformasi Laplace dari $\delta(t-a)$ seharusnya 
merupakan transformasi Laplace dari turunan fungsi Heaviside
$u(t-a)$, jika fungsi Heaviside memiliki turunan.
Pertama,
\begin{equation*}
{\mathcal{L}} \bigl\{ u(t-a) \bigr\} = \frac{e^{-as}}{s}.
\end{equation*}
Untuk memperoleh bentuk transformasi Laplace dari turunannya,
kita mengalikan dengan $s$ sehingga diperoleh $e^{-as}$, yang merupakan transformasi Laplace
dari $\delta(t-a)$.
Kita melihat hal yang sama dengan menggunakan pengintegralan,
\begin{equation*}
\int_0^t \delta(s-a)\,ds = u(t-a) .
\end{equation*}
Jadi, dalam arti tertentu,
\begin{equation*}
%mbxSTARTIGNORE
\text{``}
%mbxENDIGNORE
%mbxlatex \text{"}
\quad \frac{d}{dt} \Bigl[ u(t-a) \Bigr] = \delta(t-a) . \quad
%mbxSTARTIGNORE
\text{''}
%mbxENDIGNORE
%mbxlatex \text{"}
\end{equation*}
Alur penalaran ini memungkinkan kita membicarakan turunan fungsi dengan
diskontinuitas loncatan.
Kita dapat memandang
turunan fungsi Heaviside $u(t-a)$ sebagai sesuatu yang takhingga
di $a$, yang tepat sama dengan pemahaman intuitif kita tentang fungsi
delta.
\end{remark}

\begin{example}
Mari kita hitung ${\mathcal{L}}^{-1} \left\{ \frac{s+1}{s} \right\}$.
Sejauh ini, kita hanya menghitung transformasi invers dari
fungsi rasional sejati dalam peubah $s$.
Artinya, derajat pembilang lebih rendah daripada derajat penyebut.
Hal itu tidak berlaku bagi $\frac{s+1}{s}$.
Kita dapat menggunakan fungsi delta untuk menghitung,
\begin{equation*}
{\mathcal{L}}^{-1} \left\{ \frac{s+1}{s} \right\}
=
{\mathcal{L}}^{-1} \left\{ 1 + \frac{1}{s} \right\}
=
{\mathcal{L}}^{-1} \{ 1 \}
+
{\mathcal{L}}^{-1} \left\{ \frac{1}{s} \right\}
=
\delta(t) + 1 .
\end{equation*}
Objek yang dihasilkan adalah fungsi
tergeneralisasi dan hanya bermakna jika diletakkan di bawah integral.
\end{example}

\subsection{Respons impuls}

Seperti yang telah kita katakan, dalam persamaan diferensial
$L x = f(t)$,
kita memandang $f(t)$ sebagai masukan dan $x(t)$ sebagai keluaran.
Kita memandang fungsi delta sebagai impuls, sehingga
untuk mencari respons terhadap impuls, kita menggunakan
fungsi delta sebagai pengganti $f(t)$.
Solusi bagi
\begin{equation*}
L x = \delta(t)
\end{equation*}
disebut
\emph{\myindex{respons impuls}}.

\begin{example}
Selesaikan (carilah respons impuls)
\begin{equation} \label{eq:lteximpulseresp}
x'' + \omega_0^2 x = \delta(t) , \quad x(0) = 0, \quad x'(0) = 0 .
\end{equation}

Terapkan transformasi Laplace pada persamaan, dan nyatakan
transformasi $x(t)$ dengan $X(s)$:
\begin{equation*}
s^2 X(s) + \omega_0^2 X(s) = 1 ,
\qquad \text{sehingga} \qquad
X(s) = \frac{1}{s^2+ \omega_0^2} .
\end{equation*}
Transformasi Laplace invers menghasilkan (untuk $t > 0$)
\begin{equation*}
x(t) = 
\frac{\sin (\omega_0 t)}{\omega_0} .
\end{equation*}
\end{example}

\begin{remark}
Pembaca yang cermat mungkin menyadari bahwa tampaknya bukan $x'(0)=0$.
Namun, sebenarnya kita ingin memandang bahwa $x(t)=0$ untuk $t \leq 0$, sehingga $x(t)$
memiliki \myquote{sudut} di $t=0$, yakni,
$x'$ memiliki diskontinuitas loncatan di sana,
dan hal itulah yang menghasilkan $\delta(t)$ ketika kita mengambil $x''$.
Lihat \remarkref{remark:deltaderivofjump}.
Syarat awalnya sesungguhnya adalah $x'(0-) = \lim_{t \uparrow 0} x'(0) = 0$.
\end{remark}

Mari kita perhatikan sesuatu tentang contoh di atas. 
Dalam \exampleref{example:undampedforcedbylaplacearbitrhs},
kita memperoleh bahwa
ketika masukannya $f(t)$, solusi bagi $Lx = f(t)$
diberikan oleh
\begin{equation*}
x(t) = 
\int_0^t
f(\tau) 
\frac{\sin \bigl( \omega_0 (t-\tau) \bigr)}{\omega_0} \, d\tau .
\end{equation*}
Artinya, \emph{solusi untuk sembarang masukan diberikan sebagai
konvolusi dengan respons impuls}.  Mari kita lihat alasannya.
Untuk fungsi $f(t)$ dan $h(t)$ yang bernilai nol untuk $t$ negatif,
\begin{equation*}
(f * h)''(t) =
\frac{d^2}{dt^2}\left[
\int_0^t
f(\tau) 
h(t-\tau) \, d\tau \right]
=
\int_0^t
f(\tau) 
h''(t-\tau) \, d\tau
= (f * h'')(t) .
\end{equation*}
Kita mendiferensialkan dua kali di bawah
integral\footnote{%
Anda mungkin khawatir karena $h$ kita
sebenarnya tidak terdiferensialkan dua kali; bahkan, $h'(0)$
tidak ada sebagai turunan biasa, dan $h''$ melibatkan
fungsi $\delta$.
Untuk perhitungan ini, Anda dapat berpura-pura bahwa $h$ terdiferensialkan dua kali
dan menerapkan aturan Leibniz standar
(fakta bahwa $h(t)=0$ untuk $t < 0$ penting), atau
memperhatikan bahwa untuk fungsi semacam itu,
integralnya sebenarnya diambil pada
$(-\infty,\infty)$ alih-alih pada $[0,t]$.}.
Andaikan $h(t)$ adalah respons impuls (solusi bagi $h''+\omega_0^2 h = \delta(t)$).
Jika kita mengonvolusikan seluruh persamaan \eqref{eq:lteximpulseresp},
ruas kirinya menjadi
\begin{equation*}
f * (h'' + \omega_0^2 h) =
(f * h'') + \omega_0^2 (f * h) =
(f * h)'' + \omega_0^2 (f * h) .
\end{equation*}
Ruas kanannya menjadi
\begin{equation*}
(f * \delta)(t) = f(t).
\end{equation*}
Oleh karena itu, jika $h$ adalah respons impuls,
maka $x(t) = (f * h)(t)$ adalah solusi bagi
\begin{equation*}
x'' + \omega_0^2 x = f(t) .
\end{equation*}

Prosedur tersebut berlaku untuk setiap persamaan linear berkoefisien konstan
$Lx = f(t)$.
Jika Anda mencari respons impuls $h$
(solusi bagi $Lh = \delta(t)$),
Anda juga mengetahui cara memperoleh keluaran $x(t)$ untuk sembarang masukan $f(t)$.
Kita cukup mengonvolusikan, $x(t) = (f * h)(t)$.
Seperti yang mungkin telah Anda perhatikan dalam contoh,
respons impuls sebenarnya hanyalah transformasi Laplace invers
dari fungsi transfer, yakni,
$h(t) = {\mathcal{L}}^{-1} \bigl\{ H(s) \bigr\}$.

\subsection{Pembengkokan balok tiga titik}
\index{pembengkokan balok tiga titik}

Contoh penerapan yang sangat berbeda
yang memuat fungsi delta adalah representasi
beban titik pada balok baja.  Perhatikan suatu balok
berpanjang $L$ yang bertumpu pada dua tumpuan sederhana di kedua ujungnya.  Misalkan $x$ menyatakan
posisi pada balok, dan misalkan $y(x)$ menyatakan lendutan balok dalam
arah vertikal.  Lendutan $y(x)$ memenuhi
\emph{\myindex{persamaan Euler--Bernoulli}}%
\footnote{Dinamai menurut para matematikawan Swiss
\href{https://en.wikipedia.org/wiki/Jacob_Bernoulli}{Jacob Bernoulli}
(1654--1705),
\href{https://en.wikipedia.org/wiki/Daniel_Bernoulli}{Daniel Bernoulli}
(1700--1782), keponakan Jacob,
dan
\href{https://en.wikipedia.org/wiki/Euler}{Leonhard Paul Euler}
(1707--1783).},
\begin{equation*}
EI \frac{d^4 y}{dx^4} = F(x) ,
\end{equation*}
dengan $E$ dan $I$ adalah konstanta\footnote{$E$ adalah modulus elastisitas dan $I$
adalah momen kedua luas.  Kita tidak perlu membahas perinciannya dan cukup
memandang keduanya sebagai konstanta tertentu yang diberikan.}, dan
$F(x)$ adalah gaya yang diterapkan per satuan panjang pada posisi $x$.  Situasi
yang kita minati adalah ketika gaya diterapkan pada satu titik seperti dalam
\figurevref{lt:beambendingfig}.

\begin{myfig}
\capstart
\inputpdft{beam-bending}{%
Diagram sebuah balok di sepanjang sumbu x yang ditahan pada titik-titik tetap di
x sama dengan 0 dan di suatu x positif. Balok melendut ke bawah
oleh beban yang diberi label F kali delta kali kuantitas x dikurangi a
di salah satu titik tengah.}
\caption{Pembengkokan tiga titik.\label{lt:beambendingfig}}
\end{myfig}

Persamaannya menjadi
\begin{equation*}
EI \frac{d^4 y}{dx^4} = -F \delta(x-a) ,
\end{equation*}
dengan $x=a$ titik tempat massa diterapkan.  Konstanta $F$ adalah gaya
yang diterapkan dan tanda minus menunjukkan bahwa gaya mengarah ke bawah, yakni, dalam
arah $y$ negatif.  Titik-titik ujung
balok memenuhi syarat
\begin{align*}
& y(0) = 0, \qquad y''(0) = 0, \\
& y(L) = 0, \qquad y''(L) = 0.
\end{align*}
Lihat \sectionref{sec:appeig} untuk informasi lebih lanjut tentang syarat
titik ujung yang diterapkan pada balok.


\begin{example} \label{lt:examplebeam}
Andaikan panjang balok adalah 2, dan $EI=1$ demi
kesederhanaan.  Andaikan pula bahwa gaya $F=1$ diterapkan di $x=1$.
Artinya, kita memiliki persamaan
\begin{equation*}
\frac{d^4 y}{dx^4} = -\delta(x-1) ,
\end{equation*}
dan syarat titik ujungnya adalah
\begin{equation*}
y(0) = 0, \qquad y''(0) = 0, \qquad
y(2) = 0, \qquad y''(2) = 0.
\end{equation*}

Kita dapat mengintegralkan, tetapi menggunakan transformasi Laplace
bahkan lebih mudah.
Kita menerapkan transformasi
dalam peubah $x$, bukan peubah $t$.  Sekali lagi, kita nyatakan
transformasi $y(x)$ dengan $Y(s)$.
\begin{equation*}
s^4Y(s)-s^3y(0)-s^2y'(0)-sy''(0)-y'''(0)
= -e^{-s}.
\end{equation*}
Kita memperhatikan bahwa $y(0) = 0$ dan $y''(0) = 0$.  Misalkan
$C_1 = y'(0)$ dan $C_2=y'''(0)$.
Kita selesaikan untuk $Y(s)$,
\begin{equation*}
Y(s) = \frac{-e^{-s}}{s^4} + \frac{C_1}{s^2}+ \frac{C_2}{s^4} .
\end{equation*}
Kita mengambil transformasi Laplace invers dengan menggunakan 
sifat pergeseran kedua \eqref{ltd:sseq} untuk suku
pertama:
\begin{equation*}
y(x) = \frac{-{(x-1)}^3}{6} u(x-1) + C_1 x + \frac{C_2}{6} x^3 .
\end{equation*}
Kita masih perlu menerapkan dua syarat titik ujung.  Karena syarat tersebut
berada di $x=2$, kita cukup mengganti $u(x-1) = 1$ ketika mengambil
turunannya.  Oleh karena itu,
\begin{equation*}
0 = y(2) = \frac{-{(2-1)}^3}{6} + C_1 (2) + \frac{C_2}{6} 2^3 =
\frac{-1}{6} + 2 C_1 + \frac{4}{3} C_2 ,
\end{equation*}
dan
\begin{equation*}
0 = y''(2) = \frac{-3\cdot 2 \cdot (2-1)}{6} + \frac{C_2}{6} 3\cdot 2 \cdot 2
 = -1 + 2 C_2 .
\end{equation*}
Jadi, $C_2 = \frac{1}{2}$.
Dengan menyelesaikan untuk $C_1$ menggunakan persamaan pertama, kita memperoleh
$C_1 = \frac{-1}{4}$.  Solusi kita untuk lendutan balok adalah
\begin{equation*}
y(x) = \frac{-{(x-1)}^3}{6} u(x-1) - \frac{x}{4} + \frac{x^3}{12} .
\end{equation*}
\end{example}

\subsection{Latihan}

\begin{exercise}
Selesaikan (carilah respons impuls)
$x'' + x' + x = \delta(t)$, $x(0) = 0$, $x'(0)=0$.
\end{exercise}

\begin{exercise}
Selesaikan (carilah respons impuls)
$x'' + 2 x' + x = \delta(t)$, $x(0) = 0$, $x'(0)=0$.
\end{exercise}

\begin{exercise}
Suatu pulsa dapat datang kemudian dan dapat lebih besar.
Selesaikan 
$x'' + 4 x = 4\delta(t-1)$, $x(0) = 0$, $x'(0)=0$.
\end{exercise}

\begin{exercise}
Andaikan $f(t)$ dan $g(t)$ adalah fungsi terdiferensialkan
(dan turunannya kontinu),
serta andaikan $f(t) = g(t) = 0$ untuk semua $t \leq 0$.  Tunjukkan bahwa
\begin{equation*}
(f * g)'(t) = (f' * g)(t) = (f * g')(t) .
\end{equation*}
\end{exercise}

\begin{exercise}
Andaikan $L x = \delta(t)$, $x(0) = 0$, $x'(0) = 0$, memiliki solusi
$x(t) = t e^{-t}$ untuk $t > 0$.  Carilah solusi bagi
$Lx = t^2$, $x(0) = 0$, $x'(0) = 0$ untuk $t > 0$.
\end{exercise}

\begin{exercise}
Hitung
${\mathcal{L}}^{-1} \left\{ \frac{s^2+s+1}{s^2} \right\}$.
\end{exercise}

\begin{exercise}[menantang]
Selesaikan \exampleref{lt:examplebeam} dengan mengintegralkan 4 kali dalam peubah $x$.
\end{exercise}

\begin{exercise}
Andaikan kita memiliki balok berpanjang $1$ yang ditumpu sederhana pada kedua ujungnya dan
andaikan gaya $F=1$ diterapkan di $x=\nicefrac{3}{4}$ dalam arah
ke bawah.  Andaikan $EI=1$ demi kesederhanaan.  Carilah lendutan balok
$y(x)$.
\end{exercise}

\setcounter{exercise}{100}

\begin{exercise}
Selesaikan (carilah respons impuls)
$x'' = \delta(t)$, $x(0) = 0$, $x'(0)=0$.
\end{exercise}
\exsol{%
$x(t) = t$
}

\begin{exercise}
Selesaikan (carilah respons impuls)
$x' + a x = \delta(t)$, $x(0) = 0$.
\end{exercise}
\exsol{%
$x(t) = e^{-at}$
}
%$sX+aX = 1$
%$X = \frac{1}{s+a}$

\begin{exercise}
Andaikan $L x = \delta(t)$, $x(0) = 0$, $x'(0) = 0$, memiliki solusi
$x(t) = e^t \sin(t)$ untuk $t > 0$.  Carilah (dalam bentuk tertutup) solusi bagi
$Lx = e^t$, $x(0) = 0$, $x'(0) = 0$ untuk $t > 0$.
\end{exercise}
\exsol{%
$x(t) = (e^t) * (e^t \sin(t) ) = e^t (1-\cos(t))$
}

\begin{exercise}
Hitung
${\mathcal{L}}^{-1} \left\{ \frac{s^2}{s^2+1} \right\}$.
\end{exercise}
\exsol{%
$\delta(t) - \sin(t)$
}
%1 - 1/(s^2+1)

\begin{exercise}
Hitung
${\mathcal{L}}^{-1} \left\{ \frac{3 s^2 e^{-s} + 2}{s^2} \right\}$.
\end{exercise}
\exsol{%
$3 \delta(t-1) + 2 t$
}

%%%%%%%%%%%%%%%%%%%%%%%%%%%%%%%%%%%%%%%%%%%%%%%%%%%%%%%%%%%%%%%%%%%%%%%%%%%%%%

\sectionnewpage
\section{Menyelesaikan PDE dengan transformasi Laplace}
\label{laplacepde:section}

\sectionnotes{1--1.5 kuliah, dapat dilewati}

Transformasi Laplace berasal dari keluarga transformasi yang sama dengan
deret Fourier\footnote{Terdapat pula transformasi Fourier
pada garis real yang agak menyerupai transformasi Laplace.},
yang kita gunakan dalam \chapterref{FS:chapter} untuk menyelesaikan persamaan diferensial
parsial (PDE).
Oleh karena itu, tidak mengherankan bahwa kita juga dapat menyelesaikan PDE dengan transformasi
Laplace.

Diberikan suatu PDE dalam dua peubah bebas $x$ dan $t$, kita menggunakan
transformasi Laplace pada salah satu peubah (mengambil transformasi
dari segala sesuatu yang tampak), dan turunan terhadap peubah itu menjadi
perkalian dengan peubah hasil transformasi $s$.
PDE tersebut menjadi ODE\@, yang kemudian kita selesaikan.
Setelah itu, kita menginverskan transformasi untuk mencari solusi bagi masalah
semula.
Prosedur ini paling baik dilihat melalui contoh.

\begin{example}
Perhatikan PDE orde pertama
\begin{equation*}
y_t = - \alpha y_x, \qquad \text{untuk } x > 0, \enspace t > 0,
\end{equation*}
dengan syarat sisi
\begin{equation*}
y(0,t) = C, \qquad y(x,0) = 0 .
\end{equation*}
Kita akan mengandaikan bahwa $\alpha > 0$ adalah suatu konstanta.
Persamaan ini 
disebut \emph{\myindex{persamaan konveksi}} atau kadang-kadang
\emph{\myindex{persamaan transpor}}, dan persamaan ini sudah
muncul dalam \sectionref{fopde:section}, dengan syarat yang berbeda.
Lihat \figureref{lt:half-infinite-goo-river} untuk diagram susunannya.

\begin{mywrapfig}{3.1in}
\capstart
\inputpdft{half-infinite-goo-river}{%
Diagram bidang xt dengan kuadran pertama diarsir
dan diberi label y sub t sama dengan minus alfa kali y sub x,
sumbu x diberi label y sama dengan 0,
dan sumbu t diberi label y sama dengan C.}
\caption{Persamaan transpor pada setengah garis.\label{lt:half-infinite-goo-river}}
\end{mywrapfig}

Model fisik bagi persamaan ini adalah sungai berisi cairan kental,
karena kita tidak ingin ada sesuatu yang berdifusi.  Fungsi
$y$ adalah konsentrasi
suatu zat beracun\footnote{Sungainya sudah berupa cairan kental.
Kita tidak terlalu memperburuk lingkungan.}.
Peubah $x$ menyatakan posisi, dengan $x=0$
merupakan lokasi pabrik yang membuang zat beracun ke dalam
sungai.  Zat beracun mengalir ke sungai sehingga di $x=0$
konsentrasinya selalu $C$.  Kita ingin melihat apa yang terjadi setelah pabrik,
yakni pada $x > 0$.  Misalkan $t$ adalah waktu, dan andaikan
pabrik mulai beroperasi pada $t=0$, sehingga pada $t=0$ sungai hanya berisi
cairan kental murni.

Perhatikan fungsi dua peubah $y(x,t)$.
Mari kita tetapkan $x$ dan mentransformasikan peubah $t$.
Demi kemudahan, kita memperlakukan peubah hasil transformasi $s$
sebagai parameter karena tidak ada turunan terhadap $s$.
Artinya, kita menulis $Y(x)$ untuk fungsi hasil transformasi,
dan memperlakukannya sebagai fungsi dari $x$, dengan $s$ tetap sebagai parameter.
\begin{equation*}
Y(x)
= {\mathcal L} \bigl\{ y(x,t) \bigr\}
= \int_0^\infty y(x,t) e^{-st} \,dt .
\end{equation*}
Transformasi turunan terhadap $x$ hanyalah turunan dari 
fungsi hasil transformasi:
\begin{equation*}
{\mathcal L} \bigl\{ y_x(x,t) \bigr\} =
\int_0^\infty y_x(x,t) e^{-st} \,dt
=
\frac{d}{dx} \left[\int_0^\infty y(x,t) e^{-st} \,dt \right]
=
Y'(x) .
\end{equation*}
Untuk mentransformasikan turunan terhadap $t$ (peubah yang ditransformasikan),
kita menggunakan aturan dari \sectionref{transformsofders:section}:
\begin{equation*}
{\mathcal L} \bigl\{ y_t(x,t) \bigr\} 
=
sY(x) - y(x,0) .
\end{equation*}

Dalam kasus khusus kita,
$y(x,0)=0$, sehingga
${\mathcal L} \bigl\{ y_t(x,t) \bigr\} = sY(x)$.  Kita mentransformasikan persamaan
untuk memperoleh
\begin{equation*}
sY(x) = -\alpha Y'(x) .
\end{equation*}
ODE ini memerlukan syarat awal.  Syarat awal tersebut adalah
syarat sisi lain dari PDE\@, yaitu syarat yang bergantung pada $x$.  Semuanya
ditransformasikan, jadi kita juga harus mentransformasikan syarat ini
\begin{equation*}
Y(0) = 
{\mathcal L} \bigl\{ y(0,t) \bigr\} 
=
{\mathcal L} \bigl\{ C \bigr\} 
=
\frac{C}{s} .
\end{equation*}

Kita menyelesaikan masalah ODE $sY(x) = -\alpha Y'(x)$, $Y(0) = \frac{C}{s}$, untuk memperoleh
\begin{equation*}
Y(x) = \frac{C}{s} e^{-\frac{s}{\alpha} x} .
\end{equation*}
Kita belum selesai.
Kita memiliki $Y(x)$, tetapi yang sebenarnya kita inginkan adalah $y(x,t)$.
Kita mentransformasikan kembali peubah $s$ ke $t$.
Misalkan 
\begin{equation*}
u(t) = \begin{cases} 0 & \text{jika $t < 0$}, \\ 1 & \text{selain itu}
\end{cases}
\end{equation*}
adalah fungsi Heaviside.
Karena
\begin{equation*}
{\mathcal L} \bigl\{ u(t-a) \bigr\} =
\int_0^\infty u(t-a) \, e^{-st} \,dt
=
\int_a^\infty e^{-st} \,dt
=
\frac{e^{-as}}{s} ,
\end{equation*}
maka
\begin{equation*}
y(x,t) = {\mathcal L}^{-1} \left\{
\frac{C}{s} e^{-\frac{s}{\alpha} x}
\right\}
=
C u\bigl(t-\nicefrac{x}{\alpha}\bigr) .
\end{equation*}
Dengan kata lain,
\begin{equation*}
y(x,t) =
\begin{cases}
0 & \text{jika $t < \nicefrac{x}{\alpha}$}, \\
C & \text{selain itu.}
\end{cases}
\end{equation*}
Lihat \figurevref{lt:half-infinite-goo-river-wavefront} untuk diagram
solusi ini.  Garis berkemiringan $\nicefrac{1}{\alpha}$ menunjukkan
muka gelombang zat beracun dalam gambar saat zat tersebut meninggalkan
pabrik.
Persamaan tersebut berusaha memindahkan syarat awal ke kanan
dengan kecepatan $\alpha$,
sehingga persamaan itu juga memindahkan syarat batas (muka gelombang).

\begin{myfig}
\capstart
\inputpdft{half-infinite-goo-river-wavefront}{%
Diagram kuadran pertama bidang xt (x dan t positif).
Sumbu x, yaitu sumbu horizontal, diberi label y sama dengan 0 dan sumbu t
diberi label y sama dengan C. Sebuah garis lurus diagonal dari titik asal ke atas
dan ke kanan diberi label `muka gelombang, kemiringan 1 per alfa'. Daerah
di atas garis ini diarsir gelap dan diberi label y sama dengan C.
Daerah di bawah garis diarsir terang dan diberi label y sama dengan 0.}
\caption{Muka gelombang zat beracun adalah garis berkemiringan
$\nicefrac{1}{\alpha}$.\label{lt:half-infinite-goo-river-wavefront}}
\end{myfig}


Ssst\ldots $y$ tidak terdiferensialkan.
Bahkan, fungsi itu tidak kontinu
(tampaknya tidak ada yang pernah menyadarinya).  Bagaimana kita dapat menyubstitusikan sesuatu yang tidak
terdiferensialkan ke dalam persamaan?
Nah, bayangkan saja suatu fungsi terdiferensialkan yang sangat, sangat dekat dengan $y$.
Atau, jika Anda mengenali turunan fungsi Heaviside sebagai fungsi
delta, semuanya juga baik-baik saja:
\begin{equation*}
y_t (x,t) = \frac{\partial}{\partial t} \left[
C  u\bigl(t-\nicefrac{x}{\alpha}\bigr)
\right]
=
C  u'\bigl(t-\nicefrac{x}{\alpha}\bigr)
=
C \delta\bigl(t-\nicefrac{x}{\alpha}\bigr)
\end{equation*}
dan
\begin{equation*}
y_x (x,t) = \frac{\partial}{\partial x} \left[
C  u\bigl(t-\nicefrac{x}{\alpha}\bigr)
\right]
=
- \frac{C}{\alpha}  u'\bigl(t-\nicefrac{x}{\alpha}\bigr)
=
- \frac{C}{\alpha} \delta\bigl(t-\nicefrac{x}{\alpha}\bigr) .
\end{equation*}
Jadi, $y_t = - \alpha y_x$.
\end{example}

Transformasi Laplace sangat baik untuk persamaan berkoefisien konstan.  Salah satu
keunggulan Laplace adalah kemudahannya menangani
syarat sisi takhomogen.
Mari kita coba contoh yang lebih rumit.

\begin{example}
Perhatikan
\begin{align*}
& y_t + y_x + y = 0, \qquad \text{untuk } x > 0, \enspace t > 0,
\\
& y(0,t) = \sin(t), \qquad y(x,0) = 0 .
\end{align*}

Sekali lagi, kita mentransformasikan $t$, dan menulis $Y(x)$ untuk
fungsi hasil transformasi.  Karena $y(x,0) = 0$, kita memperoleh
\begin{equation*}
sY(x) + Y'(x) + Y(x) = 0, \qquad 
Y(0) = \frac{1}{s^2+1} .
\end{equation*}
Solusi persamaan hasil transformasi adalah
\begin{equation*}
Y(x) =
\frac{1}{s^2+1} e^{-(s+1) x}
=
\frac{1}{s^2+1} e^{-xs}
e^{-x}
.
\end{equation*}
Dengan menggunakan sifat pergeseran kedua \eqref{ltd:sseq}
dan linearitas transformasi,
kita memperoleh solusi
\begin{equation*}
y(x,t) 
=
e^{-x}
\sin(t-x)
u(t-x) .
\end{equation*}
\end{example}

Kita juga dapat mendeteksi ketika masalahnya
\emph{\myindex{tak-terumuskan dengan baik}}
dalam arti bahwa masalah tersebut tidak memiliki solusi.
Mari kita ubah persamaannya menjadi
\begin{align*}
& -y_t + y_x = 0, \qquad \text{untuk } x > 0, \enspace t > 0,
\\
& y(0,t) = \sin(t), \qquad y(x,0) = 0 .
\end{align*}
Maka masalah tersebut
tidak memiliki solusi.  Pertama, mari kita lihat hal ini dalam bahasa
\sectionref{fopde:section}.
Kurva karakteristiknya adalah $t=-x+C$.  Jika $\tau$ adalah
koordinat karakteristik, maka kita memperoleh persamaan $y_\tau = 0$
di sepanjang kurva,
yang berarti suatu solusi konstan sepanjang kurva karakteristik.
Namun, kurva-kurva ini memotong sumbu $x$ dan sumbu $t$.
Sebagai contoh,
kurva $t=-x+1$ memotong di $(1,0)$ dan $(0,1)$.  Solusi
konstan sepanjang kurva, sehingga $y(1,0)$ seharusnya sama dengan $y(0,1)$.  Namun,
$y(1,0) = 0$ dan $y(0,1) = \sin(1) \neq 0$.
Lihat \figurevref{lt:half-infinite-ill-posed}.

\begin{myfig}
\capstart
\inputpdft{half-infinite-ill-posed}{%
Diagram bidang xt dengan kuadran pertama diarsir.
Sumbu x horizontal diberi label y sama dengan 0 dan sumbu t vertikal
diberi label y sama dengan sinus t. Titik (0,1) pada sumbu t diberi label
y dari (0,1) sama dengan sinus 1 dan titik (1,0) pada sumbu x diberi label
y dari (1,0) sama dengan 0. Kedua titik dihubungkan dengan garis berlabel
t sama dengan minus x ditambah 1
dan juga dengan `y konstan sepanjang kurva karakteristik ini'.}
\caption{Masalah yang tak-terumuskan dengan baik.\label{lt:half-infinite-ill-posed}}
\end{myfig}

Sekarang perhatikan transformasinya.  Masalah hasil transformasi adalah
\begin{equation*}
-sY(x) + Y'(x) = 0, \enspace Y(0) = \frac{1}{s^2+1} ,
\end{equation*}
dan solusinya seharusnya
\begin{equation*}
Y(x) = \frac{1}{s^2+1} e^{sx} .
\end{equation*}
Yang penting, transformasi Laplace ini tidak meluruh menuju nol di takhingga!  Artinya,
karena $x > 0$ di daerah yang diperhatikan,
\begin{equation*}
\lim_{s \to \infty}
\frac{1}{s^2+1} e^{sx}
= \infty \neq 0 .
\end{equation*}
Tampaknya kita hampir
dapat menggunakan sifat pergeseran, tetapi perhatikan bahwa pergeserannya
berada dalam arah yang salah.

\medskip

Tentu saja, kita tidak perlu membatasi diri pada persamaan orde pertama, meskipun
perhitungannya menjadi lebih rumit untuk persamaan orde lebih tinggi.

\begin{example}
Mari kita gunakan Laplace untuk masalah berikut:
\begin{align*}
& y_t = y_{xx}, \qquad 0 < x < \infty, \enspace t > 0,\\
& y_x(0,t) = f(t), \\
& y(x,0) = 0 .
\end{align*}
Masalah ini berkaitan dengan batang terisolasi setengah takhingga dengan
fluks panas tertentu di salah satu ujungnya.
Susunan ini akan muncul dalam kehidupan nyata pada kasus
batang yang cukup panjang, ketika kita tidak mempertimbangkan selang waktu yang cukup lama
sehingga solusi menyadari bahwa batang kita bukan takhingga, melainkan hanya
sangat panjang.
Bagaimanapun, situasi kehidupan nyata memberlakukan syarat lain
pada solusi kita $y$.
Sebagai contoh, kita akan mengandaikan bahwa solusinya terbatas.
Syarat keterbatasan ini akan menggantikan syarat batas di
ujung takhingga.
Keterbatasan
juga cukup untuk menjamin keberadaan transformasi Laplace dalam peubah $t$.

Transformasikan persamaan dalam peubah $t$ untuk memperoleh
\begin{equation*}
sY(x) = Y''(x) .
\end{equation*}
Solusi umum ODE ini adalah
\begin{equation*}
Y(x) = A e^{\sqrt{s}x} + B e^{-\sqrt{s} x} .
\end{equation*}
Perhatikan bahwa $A$ dan $B$ bergantung pada $s$.
Karena $y$ terbatas, $Y(x)$ terbatas untuk setiap $s > 0$ yang tetap,
sehingga agar $Y(x)$ tetap terbatas ketika $x \to \infty$, harus berlaku $A=0$.

Sekarang perhatikan syarat batas di $x=0$.
Transformasinya adalah $Y'(0) = F(s)$, sehingga $-\sqrt{s}B = F(s)$.
Dengan kata lain,
\begin{equation*}
Y(x) = F(s) \frac{-1}{\sqrt{s}} e^{-\sqrt{s} x} .
\end{equation*}
Jika kita mencari transformasi inversnya dalam tabel seperti yang ada di
\Appendixref{laplacelist:appendix} (atau menghabiskan sore hari
mengerjakan kalkulus),
kita memperoleh 
%\begin{equation*}
%{\mathcal L}^{-1}\left[e^{-\sqrt{s} x}\right] =
%\frac{x}{\sqrt{4\pi t^3}} e^{\frac{-x^2}{4t}} ,
%\end{equation*}
%or
\begin{equation*}
{\mathcal L}^{-1}\left[\frac{1}{\sqrt{s}} e^{-\sqrt{s} x}\right] =
\frac{1}{\sqrt{\pi t}} e^{\frac{-x^2}{4t}} .
\end{equation*}
Jadi,
\begin{equation*}
y(x,t) =
{\mathcal L}^{-1} \left[
F(s) \frac{-1}{\sqrt{s}} e^{-\sqrt{s} x}\right]
=
\int_0^t
f(\tau) 
\frac{-1}{\sqrt{\pi {(t-\tau)}}} e^{\frac{-x^2}{4(t-\tau)}} \, d\tau .
\end{equation*}
\end{example}

Laplace dapat menyelesaikan masalah ketika pemisahan peubah gagal.
Laplace tidak terganggu oleh ketakhomogenan, tetapi pada dasarnya hanya berguna untuk
persamaan berkoefisien konstan.

\subsection{Latihan}

\begin{exercise}
Selesaikan
\begin{align*}
& y_t + y_x = 1, \qquad 0 < x < \infty, \enspace t > 0,
\\
& y(0,t) = 1, \quad y(x,0) = 0 .
\end{align*}
\end{exercise}

\begin{exercise}
Selesaikan
\begin{align*}
& y_t + \alpha y_x = 0, \qquad 0 < x < \infty, \enspace t > 0,
\\
& y(0,t) = t, \quad y(x,0) = 0 .
\end{align*}
\end{exercise}

\begin{exercise}
Selesaikan
\begin{align*}
& y_t + 2 y_x = x+t, \qquad 0 < x < \infty, \enspace t > 0,
\\
& y(0,t) = 0, \quad y(x,0) = 0 .
\end{align*}
\end{exercise}

\begin{exercise}
Untuk $\alpha > 0$, selesaikan
\begin{align*}
& y_t + \alpha y_x + y = 0, \qquad 0 < x < \infty, \enspace t > 0,
\\
& y(0,t) = \sin(t), \quad y(x,0) = 0 .
\end{align*}
\end{exercise}

\begin{exercise}
Carilah masalah ODE yang bersesuaian untuk $Y(x)$, setelah mentransformasikan peubah $t$
tersebut
\begin{align*}
& y_{tt} + 3y_{xx} + y_{xt} + 3 y_x + y = \sin(x) + t, \qquad 0 < x < 1, \enspace t > 0,
\\
& y(0,t) = 1, \quad y(1,t) =  t, \quad y(x,0) = 1-x, \quad y_t(x,0) = 1 .
\end{align*}
Jangan selesaikan masalah tersebut.
\end{exercise}

\begin{exercise}
Tuliskan solusi bagi
\begin{align*}
& y_t = y_{xx}, \qquad 0 < x < \infty, \enspace t > 0,\\
& y_x(0,t) = e^{-t}, \quad y(x,0) = 0 ,
\end{align*}
sebagai integral tentu (konvolusi).  Andaikan $y$ terbatas.
\end{exercise}

\begin{exercise}
Gunakan transformasi Laplace terhadap $t$ untuk menyelesaikan
\begin{align*}
& y_{tt} = y_{xx}, \qquad -\infty < x < \infty, \enspace t > 0,\\
& y(x,0) = 0, \quad y_t(x,0) = \sin(x) .
\end{align*}
Andaikan $y$ terbatas.
Petunjuk: Perhatikan bahwa untuk setiap $s > 0$ yang tetap,
$e^{sx}$ membesar tanpa batas ketika $x \to +\infty$
dan
$e^{-sx}$ membesar tanpa batas ketika $x \to -\infty$.
\end{exercise}

\setcounter{exercise}{100}

\begin{exercise}
Selesaikan
\begin{align*}
& y_t + y_x = 1, \qquad 0 < x < \infty, \enspace t > 0,
\\
& y(0,t) = 0, \quad y(x,0) = 0 .
\end{align*}
\end{exercise}
\exsol{%
$y=(x-t)u(t-x)+t$
}

\begin{exercise}
Untuk $c > 0$, selesaikan
\begin{align*}
& y_t + y_x + c y = 0, \qquad 0 < x < \infty, \enspace t > 0,
\\
& y(0,t) = \sin(t), \quad y(x,0) = 0 .
\end{align*}
\end{exercise}
\exsol{%
$y=e^{-cx}\sin(t-x)u(t-x)$
}

\begin{exercise}
Carilah masalah ODE yang bersesuaian untuk $Y(x)$, setelah mentransformasikan peubah $t$
tersebut
\begin{align*}
& y_{tt} + 3y_{xx} + y = x+t, \qquad -1 < x < 1, \enspace t > 0,
\\
& y(-1,t) = 0, \quad y(1,t) = 0, \quad y(x,0) = (1-x^2) , \quad y_t(x,0) = 0.
\end{align*}
Jangan selesaikan masalah tersebut.
\end{exercise}
\exsol{%
$s^2Y(x) - s(1-x^2) + 3Y''(x) + Y(x) = \frac{x}{s} + \frac{1}{s^2}, \quad
Y(-1) = 0, \quad Y(1)=0$.
}

\begin{exercise}
Gunakan transformasi Laplace terhadap $t$ untuk menyelesaikan
\begin{align*}
& y_{tt} = y_{xx}, \qquad -\infty < x < \infty, \enspace t > 0,\\
& y(x,0) = \sin(x), \quad  y_t(x,0) = 0 .
\end{align*}
Andaikan $y$ terbatas.
Petunjuk: Perhatikan bahwa untuk setiap $s > 0$ yang tetap,
$e^{sx}$ membesar tanpa batas ketika $x \to +\infty$
dan
$e^{-sx}$ membesar tanpa batas ketika $x \to -\infty$.
\end{exercise}
\exsol{%
$y=\sin(x)\cos(t)$
}

\begin{exercise}
Gunakan transformasi Laplace terhadap $t$ untuk menyelesaikan
\begin{align*}
& y_{t} = y_{xx}, \qquad 0 < x < \infty, \enspace t > 0,\\
& y(0,t) = f(t), \quad  y(x,0) = 0 ,
\end{align*}
dengan $f(t)$ suatu fungsi.
Andaikan $y$ terbatas.
Berikan jawabannya sebagai konvolusi.
\end{exercise}
\exsol{%
$
y(x,t) =
\int_0^t
f(\tau) 
\frac{x}{2 \sqrt{\pi {(t-\tau)^3}}} e^{\frac{-x^2}{4(t-\tau)}} \, d\tau
$
}
